\documentclass[12pt,a4paper]{article}
\usepackage[scheme=plain]{ctex}
\usepackage[margin=2.5cm]{geometry}
\usepackage{indentfirst}
\usepackage{setspace}
\usepackage{fontspec}
\setCJKmainfont[BoldFont=SimHei,ItalicFont=KaiTi]{SimSun}
\newCJKfontfamily\kaifont{KaiTi}[AutoFakeBold]
\usepackage{newpxtext,newpxmath}

\usepackage{lastpage}





\pagestyle{plain}

\usepackage[nobottomtitles*]{titlesec}
\titleformat{\section}{\Large\bfseries}{}{0em}{}
\titlespacing*{\section}{0pt}{3ex}{1.5ex}

\usepackage{xcolor}
\usepackage{needspace}

\setlength{\parindent}{2em}
\setlength{\parskip}{0.4em}

\doublespacing
\usepackage{tocloft}

\begin{document}
\thispagestyle{empty}

\vspace*{6cm}

\begin{center}
{\Huge\bfseries 重返城堡}

\vspace{2.5cm}

{\large 领读人：牧远}

\vspace{1em}

{\large 2026年夏 · 知合中山营共读活动}

\end{center}
\newpage
\renewcommand{\contentsname}{目录}
\tableofcontents
\newpage
\begin{center}
{\LARGE \bfseries Day 1\\[1em]准备！踏足于权力之野}
\end{center}
\vspace{3em}
\section{前哨的第一封信}
\vspace{-0.3em}{\color{gray}\small\kaifont\noindent 田野调查，即通过勘测、询问、交谈、观察等手段开展实地考察、获取研究资料的研究方法。K 是我们这支田野调查队的先遣队员，已经提前到达了“城堡”所在的村庄。这是他发回来的第一封信。他的任务是尽可能地感受这个地方、理解这个地方，然后把他看到的、想到的、困惑的东西传回来，帮助我们这些即将出发的人做好心理上和认知上的准备。读完这封信后问问自己：K 在向我们传递怎样的思想姿态？}\vspace{0.5em}
\vspace{0.3em}

\noindent\rule{\textwidth}{0.4pt}

\vspace{0.3em}



同志们：



我到了。时间很少，说关键的。



城堡看不见。不是被什么挡住了，是它本身就不让你从远处辨认。但你往下走，走进村子，会发现这里每一件小事都指向它。老板看你时的眼神，人们沉默的方式，一个年轻人半夜把我叫醒盘问时那种既傲慢又紧张的样子——他在维护某种他自己也说不清楚的东西。还有一通夜里打到总办公厅的电话，有人接。这些就是我的全部线索。



我开始意识到一件事。这里最重要的东西不直观。它在语气里，在沉默的长度里，在人们视线交错的方式里。你得让自己完全沉浸进去感受，又随时能抽身出来看清楚自己感受到了什么。这件事比我预想的难得多。



K

\newpage
\section{卢卡奇论卡夫卡}
\vspace{-0.3em}{\color{gray}\small\kaifont\noindent 卢卡奇是 20 世纪最重要的马克思主义文学理论家。在这段选文中，他通过讨论一个看似技术性的问题，小说的细节描写，揭示了卡夫卡对现实的特殊态度，即在虚无主义的基底上忠实反映扭曲的现实。我们将要采取的田野调查姿态与卢卡奇有何共同之处？又是否要保留一些必要的差异？}\vspace{0.5em}
\vspace{0.3em}

\noindent\rule{\textwidth}{0.4pt}

\vspace{0.3em}



同样的区分也适用于描写性细节。孤立地看，描写性细节可以是现实的足够真实的反映——前提是这位作家确实有才华。但描写的顺序和结构能否构成对客观现实的恰当呈现，则取决于这位作家对现实整体的态度。因为正是这种态度决定了每个个别细节在整体结构中被赋予什么样的功能。如果处理起来缺乏批判意识，其结果可能是一种任意的自然主义，因为作家将无法区分有意义的细节和无意义的细节。我认为乔伊斯就是一个典型的例子。这再次凸显了现代主义本质上自然主义的一面。



卡夫卡的情况更加复杂。卡夫卡是极少数在细节处理上具有选择性的现代主义作家之一，而不是自然主义的。从形式上看，他的细节处理与现实主义作家并无太大不同。只有当检视他的基本承诺——那条决定细节选择与编排顺序的原则——时，差异才会浮现。对卡夫卡而言，这条原则是他对超越性力量（虚无）的信念；因此，在他那些虚无主义的寓言中，艺术的统一性被打破了。



但这个问题不能以形式主义的方式去处理。在有些伟大的现实主义作品中，直接的社会历史现实被超越，细节上的现实主义是建立在对超自然世界的信念之上的。以 E.T.A. 霍夫曼为例。在霍夫曼那里，细节上的现实主义与对现实之幽灵本质的信念并行不悖。不过，细察之下，他的艺术目标与现代主义之间的差异便显现了。霍夫曼的世界——尽管带着童话般的、幽灵化的氛围——相当准确地反映了他那个时代德国的状况：一个正从扭曲的封建绝对主义走向同样扭曲的资本主义的国家。在霍夫曼那里，超自然是一种呈现德国整体状况的手段——在一个社会条件尚不允许直接现实主义描写、甚至尚未揭示出典型模式的时代。典型化的锻造在更为发达的法国要容易得多——尽管巴尔扎克偶尔也使用霍夫曼发展出来的方法（如在《和解的梅尔莫特》中）。



卡夫卡比霍夫曼更世俗。他的幽灵属于日常的资产阶级生活；而既然这种生活本身就已经是不真实的，便不需要霍夫曼式的超自然幽灵。但世界的统一性被打碎了，因为一种本质上主观的视像被等同于现实本身。帝国资本主义世界（其后来法西斯产物的先兆）所制造的恐怖——人被降格为纯粹对象——这种恐惧，原本是一种主观体验，却成为了一个客观实体。对一种扭曲的反映，变成了扭曲的反映。尽管卡夫卡的艺术方法与其他现代主义作家有所不同，但其呈现的原则是相同的：世界是超越性虚无的寓言。在卡夫卡的后继者那里，这些差异逐渐缩小或完全消失。以贝克特为例，他混合了卡夫卡式的和乔伊斯式的母题，最终产品是一种完全标准化的虚无主义现代主义。
\par\vspace{0.3em}{\footnotesize\color{gray}（Lukács, G., \textit{The Meaning of Contemporary Realism}, Merlin Press, 1963。由 DeepSeek V4 Pro 翻译。）}
\newpage
\section{观察与干预：田野中的伦理困境}
\vspace{-0.3em}{\color{gray}\small\kaifont\noindent 这段来自田野调查圣经的选文讨论了一个看似微小但至关重要的问题：研究者在田野里掏出笔记本的那一刻，发生了什么？她与被观察者的关系在那一瞬间变了——信任可能破裂，角色可能重构。然而，无涉的观察可能是不存在的，在充斥策略与抉择的田野调查中尤为如此...}\vspace{0.5em}
\vspace{0.3em}

\noindent\rule{\textwidth}{0.4pt}

\vspace{0.3em}



记速记不是简单地在笔记本或电脑上写几个字的事。由于速记往往是在研究对象附近、甚至就在他们说话做事的当下写的，因此写速记本身就是一个社会性和互动性的过程。具体而言，研究者如何记、何时记速记，可能对那些人如何看待和理解她是谁、她在做什么，产生重要的影响。关于是否该记速记、如果记的话何时记怎么记，并没有什么硬性规则。但随着在某个场景中待的时间增长，通过不断试错的好处，田野研究者可以逐渐演化出一套独特的做法，使写速记这件事适应那个场景的轮廓和限制。



\vspace{0.2em}



田野研究者还必须决定何时、何地、如何写速记。显然，低头看本子或键盘去写速记会分散研究者的注意力（即便只是一瞬间），使得对他人可能正在进行的复杂、快速而微妙的行动进行密切、持续的观察变得非常困难。但除了注意力的局限之外，记速记的决定还会对与田野中人们的关系产生重大影响。研究者努力与参与者建立紧密的联系，以便能被纳入他们生活中核心的活动。然而，在这些活动进行的过程当中，她可能体验到深深的矛盾：一方面，她可能希望通过趁话语刚说出口就记下来、趁场景正在上演就描摹下来，来保存当下的即时性；另一方面，她又可能觉得掏出一个笔记本或智能手机会毁掉这个瞬间，播下不信任的种子。参与者此时可能会把她看作一个主要兴趣在于发现他们的秘密、把他们最私密和最珍视的经历变成科学探究对象的人。



几乎所有民族志研究者在某些时刻都会感到被撕裂，一方面是自己的研究承诺，另一方面是自己真诚地融入所进入的那个世界的愿望。为了解决这些棘手的关系和道德问题，许多研究者主张：如果那些被研究的人没有完全知情和明确同意，就不应开展研究的任何环节。在这种观点看来，场景中的人必须被理解为合作者，他们与研究者积极合作，向外界讲述自己的生活和自己的文化。这种相互合作要求研究者请求许可才能撰写有关事件的内容，同时也要求尊重人们不愿揭示自己生活某些方面的意愿。



另一些田野研究者则觉得不必如此严格地被征求许可的要求所约束。一些人为这种立场辩护时坚持认为，田野研究者没有披露自己意图的特殊义务，因为所有的社会生活都包含掩饰的成分，没有人会完全暴露自己所有更深层的意图和私下的活动。另一些研究者指出，为自己而写的速记和田野笔记不会对他人造成直接伤害。这种做法当然是把那些艰难的道德和个人问题推迟到日后面对是否出版或以其他方式公开这些写作的决定时再去应对。最后，还有一些人主张向当地人隐瞒自己的研究目的，理由是所获得的信息将服务于更大的善。例如，如果研究者想要描述和公开无证工厂工人或养老院中老人的生存状况，他们就必须对控制这些场所准入权的当权者隐瞒自己的意图。
\par\vspace{0.3em}{\footnotesize\color{gray}（Emerson, Fretz \& Shaw, \textit{Writing Ethnographic Fieldnotes}, 1995。由 GLM 5.2 译。）}
\newpage
\section{田野调查程序清单}
\vspace{-0.3em}{\color{gray}\small\kaifont\noindent 《文化人类学调查——正确认识社会的方法》是国内多所高校人类学、社会学专业田野调查方法的常用参考书。以下是该书第二章“方法论（上）”的完整目录。那么，田野调查究竟是一个怎样的过程？这个过程在何种意义上是可规划或应该被规划的？}\vspace{0.5em}
\vspace{0.3em}

\noindent\rule{\textwidth}{0.4pt}

\vspace{0.3em}



\vspace{0.3em}\textbf{第二章 方法论(上)……………………………(14)}



\noindent\hspace{1em}一、调查课题的选择…………………………………(14)



\noindent\hspace{1em}二、调查方案和调查提纲的拟定……………………(15)



\noindent\hspace{1em}三、调查前的准备工作………………………………(17)



\noindent\hspace{2.5em}(一)了解地理和历史的背景材料…………………………(17)

\noindent\hspace{2.5em}(二)掌握前人调查研究成果………………………………(18)

\noindent\hspace{2.5em}(三)调查队伍的组成…………………………………………(18)

\noindent\hspace{2.5em}(四)装备、物品的购置……………………………………(19)



\noindent\hspace{1em}四、调查的开始…………………………………………(21)



\noindent\hspace{2.5em}(一)选点和安排生活………………………………………(21)

\noindent\hspace{2.5em}(二)说明来意………………………………………………(22)

\noindent\hspace{2.5em}(三)熟悉环境………………………………………………(22)

\noindent\hspace{2.5em}(四)查阅地方文献…………………………………………(23)

\noindent\hspace{2.5em}(五)人口调查………………………………………………(23)



\noindent\hspace{1em}五、如何接近人民及选择报告人……………………(24)



\noindent\hspace{1em}六、关于学习语言……………………………………(26)



\noindent\hspace{1em}七、关于参与观察……………………………………(27)



\noindent\hspace{1em}八、关于访问…………………………………………(28)



\noindent\hspace{1em}九、关于抽样…………………………………………(30)



\noindent\hspace{1em}十、材料的记录和核实………………………………(31)



\noindent\hspace{2.5em}(一)影视记录………………………………………………(32)

\noindent\hspace{2.5em}(二)声音记录………………………………………………(32)

\noindent\hspace{2.5em}(三)文字记录………………………………………………(33)

\noindent\hspace{3.5em}笔记　实录　原文　日记　图表　谱系

\noindent\hspace{2.5em}(四)关于核实材料…………………………………………(35)



\noindent\hspace{1em}十一、实物标本的搜集………………………………(35)



\noindent\hspace{1em}十二、调查成果的整理公布和存档…………………(38)
\par\vspace{0.3em}{\footnotesize\color{gray}（汪宁生，《文化人类学调查——正确认识社会的方法》，文物出版社，1996）}
\newpage
\section{K 的第一份纪实}
\vspace{-0.3em}{\color{gray}\small\kaifont\noindent K 深夜抵达村庄，被一个叫施瓦尔策的年轻人叫醒盘问，一通打到城堡总办公厅的电话确认了他的“土地测量员”身份。然而，K 真的是土地测量员吗？}\vspace{0.5em}
\vspace{0.3em}

\noindent\rule{\textwidth}{0.4pt}

\vspace{0.3em}



K到达时，已经入夜了。村子被厚厚的积雪覆盖着。城堡所在的山岗连影子也不见，浓雾和黑暗包围着它，也没有丝毫光亮让人能约略猜出那巨大城堡的方位。K久久伫立在从大路通往村子的木桥上，举目凝视着眼前似乎是空荡荡的一片夜色。

然后他去找过夜的地方；酒店里人们还都没有睡，店老板虽然无房出租，但在对这位晚客的突然到来感到极度惊讶和惶乱之余，还是愿意让他在店堂里垫一个装稻草的口袋睡觉，K同意这一安排。有几个农民还在喝啤酒，但他无意同任何人交谈，便自己去阁楼上把草袋搬下来，在炉子附近躺下了。屋里很暖和，那几个农民说话声音很低，他用困倦的双眼打量了他们一番便倒头睡下了。

然而没有多久他便被吵醒了。这时只见一个城里人装束、长着一副演员般的面孔、浓眉细眼的年轻人同店老板一起站在他身边。那些农民也都还没有走，其中几个把椅子转过来对着他们，以便看得更清楚、听得更真切些。年轻人为吵醒了K而十分客气地向他道歉，自我介绍说，他是城堡主事的儿子，然后说道：“这村子是城堡的产业，凡是在这里居住或过夜的，从某种意义上讲也算是在城堡居住或过夜，没有伯爵大人的许可，谁也不能在此居留。可是您并没有得到这样的许可，至少您并没有出示这样的证明嘛。”

K半坐起身子，捋了捋头发，仰头看着众人说道：“我迷了路，这是摸到哪个村子来了？这里是有一座城堡吗？”

“那还用问，”年轻人慢条斯理地说，这时店堂各处都有人大惑不解地冲着K摇头。“这里是伯爵大人威斯特威斯的城堡。”

“一定要得到许可才能在这儿过夜吗？”K问道，似乎想弄清刚才他听到的那些话是不是在做梦。

答话是：“是必须得到许可才行，”紧接着这个年轻人伸出胳臂，向店老板和酒客们问道：“难道竟有什么人可以不必得到许可吗？”那话音和神态里，包含着对K的强烈嘲笑。

“那么我只好现在去讨要许可了。”K打着哈欠说，一面推开被子，似乎想站起来。

“向谁去讨要？”年轻人问。

“向伯爵大人，”K答道，“恐怕没有什么别的法子了吧。”

“现在，半夜三更去向伯爵大人讨要许可？”年轻人叫道，后退了一步。

“这不行吗？”K神色泰然地说，“那么您为什么叫醒我？”

这时年轻人憋不住火了。“真是活脱脱一副盲流口吻！”他喊叫起来，“伯爵衙门的尊严必须维护！我叫醒您是想告诉您：您必须立即离开伯爵领地！”

“好了，戏做够了吧。”K用异常轻的声音说，接着又躺下去，拉过被子盖在身上。“年轻人，您太过分了点，我明天还要再考虑考虑您今天的表现的。如果一定要见证的话，酒店老板和这里的各位先生就可以作证。现在请您听清楚：我是伯爵招聘来的土地测量员，明天我的助手就要带着仪器乘车随后跟来。我因为不想失去这个踏雪觅途的好机会，所以步行前来，可惜几次迷路，才到得这样晚。现在到城堡去报到时间已经太迟，这一点我自己很清楚，用不着您来赐教。正因为这样我才勉强在这草袋上凑合过夜，而您竟然——客气点说吧——举止失礼，打搅我休息。好了，我的话说完了，晚安，诸位先生！”说到这里K翻了一个身，转向炉子去了。

“土地测量员？”他听见背后将信将疑地发问，过后又无人做声了。然而不久那位年轻人便克制住自己，用压低了的——低到可以被看作是为了照顾K的睡眠，然而又大到能让他听清楚——的声音对老板说道：“我现在就打电话去问一下。”怎么，在这个乡村小酒店里居然还有电话？唔，设备还真够齐全的。这些事，一件一件地听来也使K感到惊奇，不过总起来却又在他的意料之中。他发现，电话差不多正好摆在他的头顶上方，刚才他睡眼惺忪，没有注意到。现在，如果年轻人一定要打电话，那么他无论如何不能不打搅K的睡眠，问题只在于K让不让他打这个电话？K决定还是让他打。可是这样一来，在底下装睡便没有什么意思了，于是K又恢复了仰卧的姿势。他看见农民们怯怯地聚到一起，嘁嘁喳喳议论，看来土地测量员的到来不是一件小事。这时厨房的门早已打开，膀大腰圆的老板娘站在那里几乎把门框塞满，店老板踮着脚尖走到她跟前去告诉她刚才发生的一切。现在，电话接通，开始通话了，城堡主事已经就寝，是一位副主事——数位副主事当中某位名叫弗里茨的老爷——在那边接电话。年轻人自报姓名，说是叫施瓦尔策，接着便说他发现有一个名叫K的三十多岁的男子，衣冠不整，心安理得地在酒店里一个草袋上睡觉，枕着一个小得可怜的背囊，手边放着一根拐杖。他施瓦尔策自然觉得此人形迹可疑，而因为店主在这件事上显然失职，他施瓦尔策当然就责无旁贷地要过问此事，查明情况了。对于被叫醒盘问，对于他施瓦尔策按职责惯例作出的要将K逐出伯爵领地的警告，K的反应是很不耐烦，总之看起来火气不小，然而也许不无道理，因为他自称是伯爵大人聘来的土地测量员。在这种情况下，当然至少从例行手续上看有必要对他的话进行核实，因此，他施瓦尔策恳请弗里茨老爷在城堡总办公厅询问一下是否确有这样一位土地测量员应聘前来，并请将答复立即电话告知。

电话打完店堂里便安静下来，弗里茨在那边询问，这边在等着回话。K的神态一如既往，连头也不回，看样子一点也不急于知道结果如何，两眼茫然直视前方。施瓦尔策对事情的描述是一种不怀好意和谨小慎微的混合物，这番话给了K一种印象：城堡里连施瓦尔策这样的小人物也能很容易得到可说是外交手腕方面的训练。另外，看起来那里的人也决不偷懒；你看，总办公厅有人上夜班呢。而且显然很快就做出回答，因为这时弗里茨打电话来了。不过，好像这个回话太简短，因为施瓦尔策马上又气呼呼地把听筒挂上。“我早就说了嘛！”他大声叫道：“土地测量员，连影子也没有！这人真是个卑鄙的、信口雌黄的流浪汉，说不定还更坏呢！”此时K脑子里闪过一个念头，他想：所有的人，就是施瓦尔策、那些农民，还有店老板、老板娘，眼看就要向他猛扑过来。为了至少避一避这个凶猛的势头，他从头到脚，整个儿蜷缩到被窝里面去了。这时电话铃又响起来，而且，K觉得声音特别急促响亮。于是他又把头从被子里慢慢伸出来。虽然电话几乎不可能仍是涉及K的事情，但所有的人还是一下子突然肃静下来，施瓦尔策则回到电话机旁去。他站在那儿耐心地听完了一番较长的解释之后，低声说道：“那么是弄错了？我实在是太难为情。办公室主任亲自打来了电话？真是怪事，真是怪事。我该怎么向土地测量员先生解释才好呢？”

K听了这些话精神为之一振。这么说，城堡已经任命他为土地测量员了。这个情况一方面对他不利，因为它表明，城堡已经了解他的底细，在反复掂量了双方的力量对比之后，决定成竹在胸地开始同他较量。可是另一方面，情况又对他有利，因为在他看来这同时也证明对方低估了他，他将会有兴许比他自己起初敢于希冀的还要更多的一些自由。并且，如果以为通过这从心理战来看不能不说是高明的一着，即公开承认他的土地测量员身份，就能使他经常处于诚惶诚恐、心惊胆战的状态的话，那他们就错了；这一官方认可使他吃了一惊，但也就止此而已。

对羞怯地向他走来的施瓦尔策，K摆手示意他不必过来了；此刻人们忙不迭地恳求他赶快搬到老板的房里去住，他也拒绝了，只是接受了老板递给他的一杯催眠饮料，又从老板娘手里接过一盆水、肥皂和毛巾。现在根本用不着他开口叫人离开店堂，原来这时所有的人都把脸背过去使劲往外挤，可能是怕明天被他认出来吧。关灯之后，他终于可以休息了。他睡得很香，除了一两次被跑过的老鼠惊醒之外，一直睡到第二天早晨。
\par\vspace{0.3em}{\footnotesize\color{gray}（第一章，1-5页）}
\newpage
\section{官僚制的混乱性（Kafkaesque）}
\vspace{-0.3em}{\color{gray}\small\kaifont\noindent 这是一个印度农民的“上访”故事。研究者将这个故事归结于卡夫卡式的官僚机构的混乱运作。你认可吗？}\vspace{0.5em}
\vspace{0.3em}

\noindent\rule{\textwidth}{0.4pt}

\vspace{0.3em}



2013至2014年冬季，Shreegarh 村 15 户人家的农田被政府征用，事先没有任何通知。被征土地面积因户而异，总共约 11 英亩，户均不到 1 英亩。政府征地的目的是修路。受影响的农民均来自低种姓，主要是 kushwaha 和 vishwakarma（badhai）种姓。冬季是本德尔坎德地区的小麦种植季节，也是农民的主要收入来源。土地被征走时，小麦正在地里生长。过去曾在类似的国家干预下被迫搬迁的经历，让整个村庄弥漫着一种“似曾相识”的感觉。正如卡夫卡《审判》的主人公约瑟夫·K——“他什么也没有做错，然而，一天早晨，他被逮捕了”——这 15 户村民的土地在不经通知、不作补偿的情况下就被征走了。



一位名叫 Ramraja 的低种姓农民告诉我们，村民之间有一种普遍的看法：军队会帮他们，因为上世纪 60 年代他们村子的土地就是被政府征用去建军事设施的。村民们并不清楚应该找哪个政府部门来解决问题。30 岁的低种姓农民 Murari 说，曾有一支军队驻扎单位在村里举办过活动，当时宣布要“认领”这个村庄，当地报纸还报道了这件事。不知道该怎么办的村民们，拿着那份剪报来到了军营驻地。Murari 和另外两个村民一起去见了一位军官。他告诉我们：



\begin{quotation}\kaifont\noindent 我们给他（军官）看了那份剪报。军官让我们去法院。他说：“我们收的是 bhumidhar 地（可在农民名下登记的土地），还的也是 bhumidhar 地。现在如果你们或政府把它作为安置地来保留，那就不在我们的管辖范围之内了。”我们没法跟 sahib log（长官们）争辩。我们也去找过一个政治领袖，但没有人帮忙。\end{quotation}



bhumidhar 地的所有者拥有出售土地的法定权利。但如前所述，分配给 Shreegarh 村民的土地不可买卖。邻近村庄（如 Dosah）的地价在上涨，Shreegarh 的村民却无法从中受益。在底层语境中，质疑地方精英和政府官员被视为一种冒犯。就像《城堡》中那些村民一样，Shreegarh 的底层民众在与官员争论时，越过某一界限便会退缩。他们曾前往地区行政长官（DM）办公室要求补偿。村民们告诉我们，主管该地区行政事务的 DM 拒绝对此承担任何责任，理由是征地“依法进行”。他还说，如果村民要继续追究，可以去三百多公里外的高等法院起诉。



村民们聘请了一位来自附近村庄、在 Rajapur 执业的律师。律师向他们收取了近 10 万卢比，但案子至今连高等法院的门都没进。Murari 说，律师向村民们保证案件几个月内就能解决，并要求每个受影响的农民缴纳 3000 卢比，之后又不断索要更多。受影响的村民最近见到了这位律师，令他们沮丧的是，律师竟然要求他们重新提供土地记录的复印件和身份证件——换句话说，律师连在高等法院立案的程序都还没启动。



印度的司法系统被积压案件严重压垮，从指控确定到审判开始之间隔着极为漫长的时间。印度报业托拉斯 2014 年的报道指出，印度有 3000 万件积压案件，其中 2600 万件积压在基层法院。Kasturi（2009）估计民事案件积压时间从 9 个月到 5.4 年不等。案件的积压也导致了庭外解决倾向的增加。律师的拖延和不作为对村民们来说不是一个好兆头——他们获得正义的希望渺茫。对于 Shreegarh 的农民来说，法律就像一本字迹消褪又覆盖的羊皮卷，如同《在流放地》中那位指挥官无法辨认的图纸。协议是用英文写成的——而且是法律英文——难以理解。村里没有人精通英语去破译协议上的文字。法律语言和法律程序对村民们而言始终不可辨认。当判决以国家征地的具体形式施加到他们身上时，他们对法律知之甚少。



法律的不可接近性还体现在村民寻求帮助的各种失败尝试中。各级守门人挡住了他们的去路。地区行政部门、Lekhpal（土地记录官）、政府官员构成了阻挡村民接近正义的第一层守门人。就像卡夫卡《在法的门前》中那个守门人一样，地区行政部门把门留着——说村民们可以去法院起诉。但与此同时，土地被征了，路已经修好了，农民的损失已成既定事实。还有更强大、更高的守门人守护着法律——法院。印度的司法体系由多个层级组成。最高法院位于首都新德里。最高法院之下是各邦的高等法院。地区法院处于下一级，设在各个县府所在地。地区法院下还有审理民事和刑事案件的次级法院。高等法院——按照那位律师的说法，Shreegarh 农民的案件应该在这里提交——距离村子三百多公里。除了长途跋涉，高等法院的听证又是极其昂贵的。受影响的家庭承受不起这种法律程序的全部成本。尽管法律无处不在，对村民而言却始终遥不可及。



在卡夫卡的世界里，法律采取了自然法的形式，没有错误。国家官僚机构从同一种法律中汲取权威，也绝不承认错误。这种理性在这里以一种循环的形式展开：首先，官僚机构拒不承认自身有任何过失；其次，即便存在某种不规范，纠正的程序本身就是一个法律问题，而这个问题对穷人来说几乎不可接近。因此，官僚机构就像《城堡》里那样——变成了自我服务的机器。结果是规则变得不透明，法律对底层民众而言不可接近，就像那个在“法的门前”等待了一生的乡下人。复杂的程序和晦暗的规则，使法律和官僚程序对底层民众来说始终不可辨认。
\par\vspace{0.3em}{\footnotesize\color{gray}（Khare, A., \& Varman, R., “Kafkaesque institutions at the base of the pyramid”, \textit{Journal of Marketing Management}, 2016。由 DeepSeek V4 Pro 翻译。）}
\newpage
\begin{center}
{\LARGE \bfseries Day 2\\[1em]泊入：机遇、抉择与身份}
\end{center}
\vspace{3em}
\section{前哨的第二封信}
\vspace{-0.3em}{\color{gray}\small\kaifont\noindent K 又在村庄度过了两个白天和夜晚。两个可疑的“老助手”不请自来；城堡的第一封公函经由信使送到手上；在贵宾楼的酒吧里，他成了第十办公厅主任的情人的情人。K 在一步步接近城堡，也在一步步被城堡吸纳。}\vspace{0.5em}
\vspace{0.3em}

\noindent\rule{\textwidth}{0.4pt}

\vspace{0.3em}



同志们：



事情的发展比我预期的要快得多，也复杂得多。



先说我获得的东西。两个助手出现了，自称城堡派来的，对土地测量一无所知。我是怎么区分他们的？我用同一个名字叫他们俩。给他们任何任务，都共同承担全部责任。要么他们拆穿彼此的谎言，要么在这件事上成为同一个人。哪一种结果对我都有利。



接着信使来了。巴纳巴斯，一个举止和眼神都与旅馆里那些农民完全不同的人。他把第十办公厅主任的信交到我手上。我独自在阁楼房间里反复读了多遍。信中的措辞很微妙，既把我当成可以自由决定是否供职的人，又步步为营地把我塞进“下属”的位置。如果我接受这个位置，我猜我的行动空间会比现在大得多。不接受呢？信中没写后果。不写比写更重。



然后是弗丽达。克拉姆的情人。我透过门孔第一次看到了克拉姆本人，一个富态的、戴着夹鼻眼镜的老爷，坐在书桌后面，桌上没有文件。一束刺眼的灯光打在他身上。随后，我对弗丽达说：离开克拉姆，做我的情人。我说出口了。她答应了。



但我现在坐在这里写这些字的时候，有一个问题挥之不去。我对弗丽达说那些话，有多少是因为她这个人，有多少是因为她是克拉姆的情人？我当时以为我能把这两件事分开，现在不这么确定了。或者说，在这片田野里，情感本身就是策略的一部分，区分它们是徒劳的。



K

\newpage
\section{老助手？老助手。}
\vspace{-0.3em}{\color{gray}\small\kaifont\noindent K 回到酒店，两个自称城堡派来的“老助手”已经在等着他。他们没有仪器，不懂土地测量，而 K 照单全收。}\vspace{0.5em}
\vspace{0.3em}

\noindent\rule{\textwidth}{0.4pt}

\vspace{0.3em}



当他们快到酒店——K还记得那近处的一个拐弯——时，他大为惊异：天竟然完全黑了。他离开这里已经那么长时间了吗？可是他估算一下，大约才不过一两个小时而已，他是早上离开酒店的，现在一直还没有感到有吃饭的欲望，不久之前又还是大白天，现在可就突然漆黑一片了。“多短的白天哟！多短的白天哟！”他自言自语着，从雪橇上就势滑下来，向酒店走去。

令他十分欣慰的是，酒店老板已经站在门前的台阶上迎接他，手里还举着风灯为他照路。K走了几步又匆匆记起车夫而停了一下，黑暗中什么地方有人连连咳嗽，那就是他了。先不管他了吧，反正很快就会再见到他的。等K上了台阶，走到毕恭毕敬地向他敬礼的店老板面前时，才发现他左右两侧各站着一个人。他从老板手里接过了风灯，就着光亮细看这两个人：原来这正是他在路上遇到过的那两个被称为阿图尔和耶里米亚的男子。现在他们向他举手行军礼。这使他回想起他服役的日子，想起那些快活的岁月而笑了。“你们是什么人？”他问，同时一会儿看看这个，一会儿看看那个。“您的助手，”他们回答。“他们是您的助手，”店老板低声证实道。“怎么着？”K问，“你们就是我那两个老助手，就是我让你们随后来，我在这里等着你们的那两个人吗？”两人点头称是。“很好，”片刻之后K说道，“你们来了很好。”——“不过，”又过了片刻，K说，“你们迟到的太多了，你们太拖拉了。”——“路很远呵，”一个说。“路是很远，”K重复说，“可是你们从城堡来时我碰上你们了。”——“对，”他们说，没有更多的解释。“仪器你们都放在哪儿了？”K问道。“我们没有仪器，”他们说。“就是我交给你们保管的那些仪器，”K说。“我们没有，”他们又重复一遍。“唉，你们是怎么搞的！”K说，“你们懂不懂什么叫土地测量？”——“不懂，”他们说。“如果你们是我的老助手，那么你们就必须懂这个，”K说。他们不言语了。“好，那么来吧，”K说着便推着两人进了酒店。

之后，他们三个几乎是一声不响地坐在店堂里一张小桌旁，啤酒摆在桌上，K坐在中间，两个助手一左一右。除他们外，店里现在就只有另一张桌旁坐着几个农民，与头天晚上差不多。“同你们相处可难了，”K说，一边来回比较两人的脸庞，像他多次做过的那样，“我到底怎么分清你们两个呢？你们只有名字不同，除此以外，你们长得太像了，没有一点区别，就跟”——他停顿了一下，然后便不由自主地说下去——“除此以外你们长得就跟两条蛇一样相像。”他们微笑了。“除了名字以外人们也是很容易区别我们的，”他们分辩说。“这我相信，”K说，“我亲眼看见了，可是我呢，我只能用我的眼睛，而光用眼睛是无法分清你们两个谁是谁的。因此，我将把你们当成一个人看待，管你们都叫阿图尔，你们不是有一个叫这个名字吗？是不是你？”K问其中一个。“不是，”这人答道，“我叫耶里米亚。”——“这没什么关系，”K说，“我将管你们两个叫阿图尔。要是我说派阿图尔去哪里，那么你们两人都去，要是我让阿图尔做某件事，那么你们就两人都去做这件事。这对我虽然有一个很大的不便，就是我不能分别使用你们，然而却也有一大优点，就是对我交给你们做的任何工作，你们都共同担负全部责任。你们自己内部怎样分工我不管，但你们不能互相推诿，对我来说你们只是一个人。”他们考虑了一会儿，然后说道：“这样我们觉得太别扭了。”——“怎么能不别扭呢，”K说，“这对你们当然一定会很别扭，可就得这样定下来。”

这一阵子以来，K早看见一个农民轻手轻脚地在他们桌旁转悠，探头探脑，欲言又止，现在他终于下决心走到一个助手身后，想对他说几句悄悄话。“对不起，”K说，同时一捶桌子，站了起来，“这两位是我的助手，我们现在正在谈正事。谁都没有权利打扰我们。”——“哦，请便，哦，请便，”那农民战战兢兢地说，边说边往后退着回到和他一起的那帮人那里去了。“这一点你们两个务必注意，”K重新坐下后说道。“没有我的许可，你们不得同任何人谈话。我在这里是个外乡人。既然你们是我的老助手，那你们也就是外乡来的。因此，我们三个外乡人要齐心，伸出你们的手，让我们拉拉手吧。”两人忙不迭地向K伸出手来。“你们两个也握一握手吧，”他说，“对我要有令则行。我现在要去睡了，我劝你们也去睡一觉。今天我们耽误了一个工作日，明天必须很早开始工作。你们要设法找一辆雪橇来，我们好去城堡，六点钟必须一切准备停当在这门口待命。”——“好的，”其中一人说。但另一个却抢白说：“你说‘好的’，可你明明知道这办不到。”——“不要说了，”K说，“难道你们现在就要一个跟一个闹分裂？”但是，现在第一个助手也否认了：“他说得对，没有许可任何外人是不得进入城堡的。”——“要到哪里去申请许可呢？”——“我不知道，也许是找主事吧。”——“那我们就打电话到那儿去申请，你们两个听着：马上打电话给主事！”两人立即跑步来到电话机旁，接通了主事——瞧他们那争先恐后的劲头儿！从表面看，他们真是唯命是从、唯唯诺诺到可笑的地步——问道，K可不可以同他们一起上城堡。电话里厉声回答：“不行！”就连K坐在他桌旁也听见了。然而答话比这还要多一点，叫做：“明天不行，任何时候都不行。”——“我自己来打，”K说着站了起来。到目前为止，除了那个农民引起的干扰外，很少有人注意K和他的两个助手，然而现在他这最后一句话却引起了众人的注意。所有的人一下子跟他一齐站了起来，然后，虽然老板力图把这些人推开，他们仍在电话机旁围成一个半圆形的圈子将K包围了。他们议论纷纷，占主导地位的看法是：K将得不到任何回答。K不得不请求他们安静下来，说他并没有要求他们把他们的意见讲给他听。

听筒里传来K以往打电话时从未听到过的嗡嗡声。听起来，就像从一大片乱哄哄的孩子吵嚷声中——可这声音又不是真正的嗡嗡声，而是从远方，从很远很远的地方传来的歌唱声，就像是从这一片嗡嗡声中神奇而不可思议地逐渐幻化出一个单一的、很高的强音，它猛烈撞击着他的耳鼓，仿佛它强烈要求深深钻入人体内部而不只是接触一下那可怜的听觉器官似的。K全神贯注地等候，一句话不说，左臂支在摆电话机的小桌上，就这样聚精会神地恭候着。

他不知过了多长时间；只知道后来老板拉扯他的上衣，告诉他有一个信差来找他。“走开！”K忍不住大叫一声，也许这吼声传进听筒里去了吧，因为那边马上有人接电话了。于是他听见了下面的话：“我是奥斯瓦尔德，那边是谁呀？”听筒里叫道，这是一个严厉而高傲的声音，K觉得有点小小的语音缺陷，那声音力图超越自身，用更加咄咄逼人的凌厉气势来弥补这个缺陷。K在犹豫，没有立即报上自己的名字，在电话面前他是赤手空拳，毫无防御能力的，对方可以厉声呵斥他，可以把听筒撂下，这样一来K就无异于自己把一条也许是相当重要的通路堵死了。K的犹豫使那边的男子不耐烦起来：“打电话的谁呀？”重复喝问之声接着又补充道：“请你们那边不要老打电话来好不好？刚刚才来过一个电话嘛！”K没有理睬这句话，他突然灵机一动，报上自己身份说：“我是土地测量员先生的助手。”“什么助手？什么先生？什么土地测量员？”K想起了昨天的电话。“请您问问弗里茨吧，”他简短地说。使他吃惊的是这话居然奏效了，比这更使他惊讶的是那边办事口径完全一致。答话是：“我知道了。那个永远扯不清的土地测量员。是了，是了。还有什么？哪个助手？”“约瑟夫，”K说。他对那些农民在他身后嘀嘀咕咕有些不快；显然他们不同意他没有报上真实身份这一做法，但K没有功夫去考虑他们，因为应付电话就够使他费神的了。“约瑟夫？”那边反问道。“两位助手不是叫”——说到这里那边稍稍停了一下，显然是在让谁告诉名字——“阿图尔和耶里米亚吗？”“他们是新来的助手，”K说。“不，他们是老助手。”——“他们是新的助手，我才是老助手，是今天比土地测量员先生晚一步后到这里的。”——“不对！”那边大叫起来了。“那么我是谁呢？”K问道，一直保持着冷静。这样过了一会儿，听筒里又响起同一个声音，带着同样的语音缺陷，可同时又的确像是另一个更低沉、更令人肃然起敬的声音：“你是老助手。”

K过于专注地琢磨那声音的韵味，以至几乎没有听见对方的问话：“你想干什么？”〔2〕现在他真恨不得立刻放下听筒。他已不再期望能从这次谈话中得到点什么。他只是硬着头皮匆匆再问了一句：“什么时候我的主人可以到城堡来？”——“什么时候都不行，”就是回答。“知道了，”K说着便挂上了听筒。
\par\vspace{0.3em}{\footnotesize\color{gray}（第二章）}
\newpage
\section{一封危险的公函}
\vspace{-0.3em}{\color{gray}\small\kaifont\noindent 信使巴纳巴斯，一个看似超脱于村庄的年轻人，递上了第十办公厅主任署名的公函。信中的措辞模棱两可：K 是做微末的下属，还是主任身边的透明人？}\vspace{0.5em}
\vspace{0.3em}

\noindent\rule{\textwidth}{0.4pt}

\vspace{0.3em}



这时他身后那些农民早已挤到他跟前来了。两个助手一边不断斜眼看K，一边不住地挡住农民不让他们靠近。然而他们只不过是做做样子罢了。农民们实际上也已经对谈话结果感到满意而逐渐后退了。正在这时，从他们后面，一个人大步流星，把农民向两边分开走了过来，他向K鞠了一躬，递交给他一封信。K手里拿着信，两眼注视着来人，觉得这人眼下比别的事情更加重要。他和两个助手有许多相似之处，同他们一样身材瘦长，一样穿着紧身衣，也同他们一样敏捷灵巧，可又与他们完全不同。要是我K的助手是这个人多好！这人某些地方使K想起在盖尔斯泰克那里见过的那个怀抱婴儿的女人。他的着装几乎全是白色，那套衣服大概不是丝绸的，它是很普通的冬衣，但却具有丝绸衣服的柔和和庄重。他的脸明亮而开朗，眼睛非常大。他的微笑有出奇的感染力；他以手拂面，那神气似乎想赶走这微笑，然而却不成功。“你是谁？”K问道。“我叫巴纳巴斯，”他说。“我是个信使。”说话时他那两片嘴唇一张一闭，既显出阳刚之气，又饱含柔和之美。“你喜欢这里吗？”K问道，同时指了指那伙农民，到现在他一直还保持着对他们的兴趣，这些人有着一张张简直就是受苦受难的脸——他们的颅骨似乎被人打扁了，面部线条则是挨打后的痛苦表情勾勒出来的——和厚大的翻嘴唇，他们张开大嘴站在一旁观看，但同时又没有观看，因为有时他们的目光游移不定，茫然若失地盯住某件无关紧要的东西瞅一阵，然后才又返回到他和巴纳巴斯身上来；除这些农民外，K又指了指他的两个助手，这两人现在紧紧地手拉着手，脸贴着脸，面带微笑，无法确定这笑是卑逊的呢还是轻蔑的，他把他们全都指给他看，仿佛在介绍特殊境遇强加于他的一批随从，并期望着——这期望里有一种一见如故的心情，K很希望同这个人谈谈心里话——巴纳巴斯时时把他和这些人区别开来。但是巴纳巴斯根本就不回答——当然不含任何恶意，这是看得出来的——他的问题，恰似一个训练有素的仆人在听到主人仅仅表面上针对他而说的话时露出的表情那样，顺从地默默听着，只是遵循所提问题的意思左顾右盼，然后举手向农民中的熟人打招呼，又同两个助手交谈了几句，言谈、举止、神态是自由而自主的，一点没把自己同他们混同起来。K转眼——他问话碰了个软钉子，但并不觉难堪——看看手中的信并拆开了它。这封信的全文是：“非常尊敬的先生！如您所知，您已被聘任前来为大人供职。您的直接上司是村长，他还将告知您有关您的工作及薪俸的一切细节，而您也将负责向他汇报工作。尽管如此，我本人还是要在百忙中兼顾您的。递送此信的巴纳巴斯将不时向您询问以了解您的愿望，然后及时禀告我。您将发现我会尽一切可能为您提供方便。使我的部下心满意足，实为我所期盼。”签署的名字很潦草，看不清楚，但在签署旁边印有：“第十办公厅主任”字样。“你等一等！”K对正向自己鞠躬的巴纳巴斯说，然后他叫老板要一个房间用用，说他想一个人再好好看看这封信。同时他又想到，不论自己对巴纳巴斯有多少好感，他毕竟只是个信差，于是吩咐拿些啤酒给他。K留意观察他对此的反应，看到他显然很高兴地接受了这一赏赐，立刻喝起来。然后K便随老板而去。在酒店这所小房子里，只能做到为K准备一间阁楼斗室，而即使这一点也不那么容易，因为必须把原来在这房间里睡觉的两个女仆安排到别处去住。其实，也只不过是让她们出去而已，房间本身看来原封未动：唯一的一张床上，没有床单、枕套等物，只有几块垫子和一床粗毛毯子，还保持着昨夜使用过以后的原状。墙上挂着几张圣徒画像和几张士兵照片。屋子里甚至没有开窗换换空气，显然是希望新来的客人不要在这里长住，所以不做任何能吸引他留下的事。但K挺能凑合，怎么都行，他披上毯子，坐到桌前，开始借着一支蜡烛的光亮再次读起那封信来。

这封信的内容前后并不一致，有几处像在同一个享有充分自由的人谈话，他自己的意愿受到尊重，如信头，还有涉及他的愿望的那句话都是如此。然而信中也有一些地方或隐晦或明显地把他当成一个小小的、从办公厅主任坐椅上看去小得几乎看不见的下属，主任必须作出很大的努力，才能“在百忙中兼顾”他一下，他的上司只是村长，K甚至还有向这位村长汇报工作的责任，他的唯一同事恐怕就只有那个村警了。这无疑都是一些矛盾，它们过于明显，使人觉得一定是有意这样措词的。K几乎不敢在这样的官府面前不知天高地厚地认为这里可能有犹豫不决的因素在起作用，相反，他觉得信中是明白无误地给他提出了两种可能的选择，完全让他自己决定如何对待信中作出的安排：是想做一名同城堡有着说起来倒是很光彩、然而却只是一种表面上的联系的农村工作人员呢，还是做一名仅仅是表面上的农村工作人员，而事实上却由巴纳巴斯带的信息来决定他和城堡的全部工作关系？K在选择上毫不犹豫，而且即便他没有到这里之后取得的那些经验，他也是不会犹豫的。只有作为一名农村工作人员，离城堡老爷们愈远愈好，他才有可能在城堡达到一定的目的。村里这些人，现时对他还抱着很大的不信任，而一旦他不说成为他们的朋友，就算只成为同他们一样的村民，他们就会开口说话，而一旦他同盖尔斯泰克或拉塞曼没有任何区别——这一点必须迅速做到，这是关键的一步，那么，在他面前所有的道路就一定会一下子畅通无阻，而如果他仅仅寄希望于山上那些老爷们及其恩典，他就不仅会永远被阻挡在这些大路之外，而且恐怕连这些路的影子也永世看不到。当然，存在着一种危险，信中反复强调了这种危险，在表述它时字里行间流露出某种喜悦，似乎它是摆脱不掉的了。这危险就是他的下属地位：供职、上司、工作、薪俸、汇报、部下，这封信中充斥着这些字词，就是在谈到别的问题，比较涉及个人的问题时，也是从这一角度出发的。如果K愿意，那么他马上可以成为这个部下，但那就是极为严肃的事，没有丝毫另就的可能。K知道人家并没有拿要采取任何真正的强迫手段来威胁他，他不怕强迫，尤其在这个地方更不怕，但是他害怕这个使人泄气的环境对他产生的强大压力，害怕对各种失望逐渐习以为常的那样一种习惯势力，害怕无时无刻不在向他袭来的各种潜移默化的影响形成的强大压力，然而无论怎样害怕，他必须冒着这种危险勇敢地投入战斗。信上不是也不讳言，如果真需要靠斗争来一决雌雄，那么K是有胆量先发制人的吗；话说得很含蓄，唯有一颗焦灼不安的心——是焦灼不安的心，而不是深感歉疚的心——才能觉出这里的弦外之音，正是“如您所知”那几个关于他被录用为大人供职的字眼，暗含着这层意思。K只是报了到，但自那时起他便如信中所说，已经知道被录用了！

K从墙上取下一张画，把信挂在那个钉子上；既然他将在这间屋子里住，那么信就理应挂在这里。
\par\vspace{0.3em}{\footnotesize\color{gray}（第二章）}
\newpage
\section{双向奔赴的爱情？}
\vspace{-0.3em}{\color{gray}\small\kaifont\noindent 在巴纳巴斯的姐妹奥尔嘉的带领下，K 在贵宾楼酒吧遇见了第十办公厅主任克拉姆的情人弗丽达。他选择向通往城堡权力核心的最短路径发起冲击。他成功了吗？}\vspace{0.5em}
\vspace{0.3em}

\noindent\rule{\textwidth}{0.4pt}

\vspace{0.3em}



在酒吧——一间很大的、当中完全空着的屋子——里，立着一箱箱啤酒，农民们有的在酒桶边，有的在酒桶上靠墙坐着，但他们的模样与K住的那家酒店里那些农民完全不同。他们衣着整洁一些，都穿一色浅土黄粗布，上衣鼓起，裤子紧身。这些男人全是矮个儿，乍一看长相都差不多，都是扁平脸庞，额头、颧骨、下巴突出，却又有多肉的面颊。他们一概少言寡语，几乎纹丝不动地坐着，只用目光尾随新来者，然而眼神又是迟滞、冷漠的。尽管如此，因为人多，又非常安静，他们仍然引起了K一定的注意。他再次拉起了奥尔嘉的胳臂，算是告诉众人他是怎么来到这里的。这酒吧的一角有一个男人站了起来，这是奥尔嘉的一个熟人，他想到这边来，可是K立刻挎起她的胳臂推着她朝另一个方向走了。这一点除她本人外谁都没觉察到，她微微一笑，斜睨他一眼，算是默许他这个行动了。

斟啤酒的是一个名叫弗丽达的少女。这是个头发金黄、脸颊瘦削、有着一双忧伤眼睛、不引人注目的小个子姑娘，可是她的目光却令人吃惊，那是一种特别高傲的目光。当它落在K身上时，K觉得它似乎已经把所有与他有关的事情统统解决了，他本人现在还一点不知道有这些事情，然而这目光却使他坚信确有其事。K一直不停地从侧面看着弗丽达，就是当她已经在同奥尔嘉说话时也还在打量她。看样子，奥尔嘉和弗丽达可不大像是朋友，她们只是冷冷地交谈了几句。K想帮她们缓和一下气氛，所以突然动问道：“您认识克拉姆先生吗？”奥尔嘉忍不住笑起来。“你笑什么？”K生气地问。“我并没有笑呵，”嘴上这么说，实际上仍不停地笑。“奥尔嘉这姑娘还真够孩子气的，”K说，同时靠着柜台把身子使劲向下弯，以便再次把弗丽达的注意力吸引到自己身上来。但她仍低垂眼皮，轻声说道：“您想见克拉姆先生吗？”K说请求一见。她便指着一扇正好就在她自己左手边的门。“这门上有个小孔，您可以从这儿往里面看。”——“那么这儿的这些人呢？”K问道。她撅了撅下嘴唇，用她那异常柔软的手把K拽到门边。把眼睛贴在那个显然是为了窥视而钻出的小孔上，他几乎可以将隔壁房间一览无余。

屋子中央，一张书桌旁有一把舒适的圆形靠椅，椅子上坐着的正是克拉姆先生，他面前挂着一盏亮得刺眼的电灯。这是一位中等身材、颇为富态、看来一定行动不甚方便的老爷。他的脸还算光溜，但面部肌肉却已经由于年龄的分量而有些下垂了。黑色的小胡子长长地伸向两边。一副歪戴着的不断反光的夹鼻眼镜挡住了他的眼睛。如果克拉姆先生完全伏案工作，那么K就只能看见他的侧面；但因为他现在几乎是面对K坐着，所以K可以看清他的整个面庞。克拉姆左肘支在桌上，拿着一支弗吉尼亚雪茄的右手搁在膝盖上。桌上摆着一个啤酒杯；由于桌子的边条较高，K不能看清桌上是否有任何文件之类，但他觉得桌面似乎是空的。为保险起见，他请弗丽达也从小孔中看一看然后告诉他究竟有还是没有。但是，因为她不久前刚到过那屋里，所以便确有把握地向他证实那桌上并没有什么文件。K问弗丽达他现在是否必须走开，可是她说如果他乐意，趴在门上看多久都行。现在K身边只有弗丽达，奥尔嘉呢，K匆匆注意到她已经找到了她的熟人，这时高高坐在一只桶上，脚不停地荡来荡去。“弗丽达，”K小声说，“您同克拉姆老爷一定很熟啰！”——“唔，是的，”她说。“是很熟。”她靠墙站在K身边，并且——这时K才注意到——用手轻轻地抚弄她那件薄薄的肉色开领衬衫，这件衣服同她那可怜的瘦小身躯很不相称。俄顷她说：“您不记得奥尔嘉刚才的笑了吗？”——“记得，这个淘气鬼，”K说。“不过，”她带着和解的口吻说，“她笑是有道理的。您刚才不是问我是不是认识克拉姆吗？可我恰好是，”——说到这里她不由自主地挺了挺胸脯，同时她那与谈话内容毫无关系的高傲目光再次越过K的头投向别处，“我恰好是他的情人呵。”——“克拉姆的情人，”K说。她点点头。“那么在您面前，”为不让他们之间的气氛过于严肃，K微笑着说，“我必须肃然起敬了。”——“不光是您，”弗丽达和蔼地说，但并不是受到了他那微笑的感染。K想到一个煞一煞她的傲气的办法，现在使了出来；他问道：“您到城堡里去过吗？”可这难不倒她，她立即答道：“没去过，可是我在这儿这个酒吧里，这不就足够了吗？”她的虚荣心显然非常强，看来此刻正想在K的身上满足一下这种虚荣心。“当然啦，”K说，“在这个酒吧里，您是很擅长做酒店老板工作的嘛。”——“是这么回事，”她说，“可我一开始时是在‘大桥’酒店当喂牲口的女佣人呀。”——“难道就用这双这么娇嫩的手喂牲口吗？”K带着疑虑的语气说，连他自己也不清楚他这话究竟只是为了说点好听的呢，还是他真的已经被她俘虏了。她的两只手确实小巧而又娇嫩；但是，管它们叫做干瘪瘦弱、索然无味也未尝不可。“当时谁也没有注意这个，”她说，“就是现在——”K疑惑地看着她。她摇摇头不想说下去。“当然，您有您的秘密，”K说，“您是不会同一个刚认识半小时，还没有机会向您介绍自己情况的人谈论这些秘密的。”可是，这话看来说得很不合适，就好像弗丽达本来处在一种对他很有利的迷迷糊糊的状态，而他竟一下子把她摇醒了。她从自己腰带上挂着的皮包里拿出一小块木头把门上小孔堵住，对K说话时显然在竭力抑制自己，不让他觉察她内心的变化：“至于说到您，那么我知道您的底细，您是土地测量员，”然后又补充道：“现在我得去干活了，”说完便走到柜台后面她的位置上，时不时有一两个酒客站起来，到柜台前让她斟酒。K还想再同她悄悄谈一次，因此从一个架子上取下一只空杯向她走去。“我只想再说一句，弗丽达小姐，”他说，“从一个喂牲口的女佣人，苦干到当上酒吧女招待，是非常了不起的，需要有超出常人的毅力，可是对于这么个不寻常的人，最终目标是不是就到此为止了呢？这是个荒谬的问题。请别笑我，弗丽达小姐，从您的眼睛里我看到的不完全是您过去的奋斗，而更多的是您将来的奋斗。但这个世界上各种阻碍是很大的，目标越大阻碍也越大，所以，希望争取到一个无权无势但也在奋斗着的小人物的帮助，便不是什么可耻的事情了。也许我们两人可以好好地单独谈谈，别有那么多双眼睛盯着我们说话。”——“我不知道您想干什么，”她说，但这次与她的意愿相反，语调中已听不到她在生活中屡次得胜产生的傲气，而只有数不清的失望引起的凄楚了。“您兴许想把我从克拉姆身边拉走吧，我的老天！”她两手合拢一击掌。“您把我看透了，”K说，显得像是被太多的不信任弄得疲惫不堪，“这一点正是我最隐秘的意图。您最好离开克拉姆，做我的情人。好了，现在我可以走了。奥尔嘉！”K叫道。“我们回家了。”奥尔嘉顺从地就势从桶上滑下，可还不能马上摆脱那帮包围她的朋友。这时弗丽达瞪了K一眼，低声说：“我什么时候可以同您谈谈？”——“我可以在这儿过夜吗？”K问。“可以，”弗丽达说。“现在就可以留下吗？”——“您先同奥尔嘉离开这里，我好设法支走眼前这帮人。过一会儿您就可以来了。”——“好吧，”K说，焦急地等着奥尔嘉。
\par\vspace{0.3em}{\footnotesize\color{gray}（第三章）}
\newpage
\section{强势田野与“真问题”}
\vspace{-0.3em}{\color{gray}\small\kaifont\noindent 一部分现代政府研究者以研究强势田野自居，强调对“真问题”的识别能力。存在一套“真问题”掩护机制，以及克服这种掩护机制的一般做法吗？}\vspace{0.5em}
\vspace{0.3em}

\noindent\rule{\textwidth}{0.4pt}

\vspace{0.3em}



（四）“强势田野”的真问题识别困境



“强势田野”意味着研究者极易陷入真问题识别困境，具体表现为三种状态：其一，受信息不对称与自身认知局限，研究者无力穿透田野的表象，难以捕捉具有研究价值的核心议题，陷入无问题的迷茫状态；其二，在表层信息供给与应付式互动的误导下，研究者易将非真实存在的议题误判为核心研究对象，陷入伪问题的认知陷阱；其三，研究者可能聚焦非核心、低学术价值的真实议题，误将边缘问题作为主要研究方向。



无问题是指研究者在“强势田野”中难以识别出具有研究价值的真实问题，陷入“看得见却看不清、接触到却摸不透”的困境。这源于以下两个方面：一方面，“强势田野”存在严格的信息筛选与层级化沟通，研究者往往被限定在公开场景与表层互动中，难以触及核心工作流程与真实运作——田野主体出于风险规避、工作保密或程序规范等考量，仅向研究者展示经过包装的标准化场景，如预先安排的观摩环节、统一口径的访谈回应、公开可查的常规资料，而那些隐藏在流程背后的矛盾冲突、政策执行弹性、主体间利益博弈等关键信息并未向研究者公开，使研究者如同“隔着玻璃看风景”，无法感知田野的真实肌理；另一方面，研究者自身的认知局限与身份隔阂增加了发现问题的难度，“强势田野”往往具备高度专业化的运作规则、独特的行业术语与隐性的行动逻辑，使研究者既难以理解田野中各类现象的深层含义，也无法判断哪些现象背后隐藏着值得探究的问题，进而只能被动接收碎片化、表面化的信息，无法将零散现象串联成有逻辑的问题线索，最终陷入“无问题可究”的迷茫状态。更为严重的是，这可能导致研究者将“无问题”误判为田野的真实状态，或将自身的认知盲区等同于田野的平稳运行，进而得出“该田野不存在值得研究的问题”的错误结论，错失挖掘核心研究议题的机会。



伪问题是指研究者在“强势田野”中识别的问题并非真实存在的问题。这源于以下两个方面：一方面，研究者对“强势田野”理解不深，进入田野如同“刘姥姥进大观园”，一切新鲜事物对其都是“问题”；另一方面，田野主体告知研究者一些老生常谈的信息，如工作的推进源于领导重视、建立领导小组、建章立制、群策群力、实现创新等，这些访谈内容实则属于“放之四海而皆准”的信息。“放之四海而皆准”意味着这些可能与研究者所提出的问题不符的内容很可能被误认为是珍贵的访谈材料。与此同时，田野主体可能提供给研究者诸多文本素材，其多以规章制度、工作总结、工作报告为主，这些材料中不乏有效的田野素材，但研究者可能难以从这些材料中感知真正有用的信息——因为许多材料大同小异，而写材料的逻辑又使文本极为凝练。有鉴于此，一些将“某些任务推进何以成功”的原因归结于上述所说建章立制、群策群力等的观点，很可能并未挖到根源，而由此去探寻某个政策执行的逻辑容易发现伪问题。



边缘问题是指研究者在“强势田野”中识别出的真实存在却不值得着重研究的问题。这源于以下两个方面：一方面，研究者聚焦“强势田野”中的非核心工作，这些工作多是各部门的通用型工作，专业化程度低意味着研究者易于上手，有时研究者还会认为这些“新事物”值得研究；另一方面，田野主体提供给研究者若干亟待解决的问题，其大多是真实存在的问题，如部门运作机制的优化。这些问题多为田野主体站在研究者角度所提出的问题：作为高级知识分子，研究者应去解决实践者解决不了的问题。显然，这些在政府部门普遍存在的问题是真实的，但可能鉴于其微观特性、技术的难解决性以及既有研究已经提供可行的方案而不能提炼为高质量的研究问题。值得注意的是，研究者往往会把边缘问题当作与“强势田野”中主体破冰的“第一桶金”，但这存在风险：如果研究者没有完成，那么田野主体可能怀疑研究者能力并减少互动意愿；如果研究者完成，那么田野主体可能持续交给研究者类似的工作，研究者没有时间进行深入的思考；更有甚者，研究者可能将这些问题当作研究问题，陷入“只见树木，不见森林”的困境。



上述三种情形的核心症结在于研究者未能穿透“强势田野”的表象，触及其实质矛盾与深层逻辑，最终导致具有学术价值的真问题缺失。破解这一困境的关键在于习得并运用地方性知识——在政府研究中，地方性知识有助于研究者把握“强势田野”的运作规则、互动惯例与隐性知识，进而突破信息壁垒、纠正认知偏差，其在问题识别、界定、正当化以及任何可能的解决方案生成过程中都扮演重要角色。
\par\vspace{0.3em}{\footnotesize\color{gray}（凌争，政府研究中的“强势田野”与地方性知识，《行政论坛》2026年第1期，第76-85页）}
\newpage
\begin{center}
{\LARGE \bfseries Day 3\\[1em]村与城：如何接近远端的存在？}
\end{center}
\vspace{3em}
\section{前哨的第三封信}
\vspace{-0.3em}{\color{gray}\small\kaifont\noindent K 在村庄继续推进他的工作。他试图靠近城堡——不止是从物理意义上靠近，更是从理解上靠近。他走到了村子的边缘，走进了一户陌生人的屋子，也走进了村公所。他看到了三种不同的东西，但他在信里没有告诉我们答案。他想问的是一个比答案更先的问题：如果你面对的是一个你看不到的东西，你该从哪里开始？}\vspace{0.5em}
\vspace{0.3em}

\noindent\rule{\textwidth}{0.4pt}

\vspace{0.3em}



同志们：



城堡在哪里？我已经从三个不同的方向试过了。结果互不相同。



先说最近的。昨天晚上我决定直接走过去。大路看起来是朝山上去的，走了很久，我确认了一件事：这条路永远不会到城堡。它在山脚下拐了弯，然后跟山平行，一直往前走，村子也不到头。不是“路走错了”那种感觉，是“路本身就是这么设计的”。



然后我进了一间屋子。屋子挺大的，有人洗澡，有人洗衣服，有个女人抱着婴儿坐在靠椅上，头也不抬，衣服泛着丝绸的光。她说是从城堡来的少女。其他人呢？不是我预想中的敌人也没有仇恨。老人给我推了一块木板让我从雪里爬出来，两个男人把我拽到门边送我出去，自始至终没有说一句“不可以”——但动作就是答案。在这间屋子里，城堡不在山顶上，它在每一个人看别人的方式里、在说话的次序里、在关门和开门的时机里。



最后我去了村公所。村长躺在床上，风湿病犯了。他讲了两件事：电话和公文。打到城堡的电话会让最下一级部门的全部电话一起响起来，但即便通了也不能当真——那是官员用来消遣的；公文在各部门之间来回传递，每个环节都有答复，但从没有裁断。



三样东西。空间上的隔绝，人际关系里的渗透，制度上的错位。它们指向的是同一个城堡吗？如果是，我该怎么把它们拼在一起？如果不是——那这个“城堡”到底是什么？



K

\newpage
\section{通往城堡的路}
\vspace{-0.3em}{\color{gray}\small\kaifont\noindent K 决定直接用双脚靠近城堡。大路看上去是朝山上去的，但走到山脚下时，路拐了弯，从此与山平行。他不是在这条路上迷了路，他是发现了这条路本身的目的。}\vspace{0.5em}
\vspace{0.3em}

\noindent\rule{\textwidth}{0.4pt}

\vspace{0.3em}



此刻，他在明澈的空气中看清了城堡的轮廓，且由于那层山上处处皆是薄薄的积雪把任何物体的形状都勾勒出来，形状益发明晰可辨了。山上的积雪看来比村里少得多。K此时在村里踏雪前进，并不比昨天在大路上省力。这里的积雪一直堆到与那些简陋小屋的窗户齐高，稍往上一点，低矮的屋顶上也沉甸甸地堆满了雪，而山上的那些建筑物，却都那么自由自在、轻松愉快地挺立着，至少从这里看上去的印象是如此吧。

站在这里，从远处看去，城堡的外观大体与K的预料相符。它既不是一座古老的骑士城堡，也不是一座新式的豪华建筑，而是一座宽阔的宫苑，其中两层楼房为数不多，倒是有许许多多鳞次栉比的低矮建筑；如果事先不知道这是一座城堡，那么可能会以为它不过是一个小镇。K只看见一座塔，它究竟是属于一所住宅还是属于一所教堂，就无法看清了。一群群乌鸦在围着它盘旋。

K心里什么别的都不想，两眼看准城堡继续朝前走。但是愈走近城堡，他就愈觉失望，那的确不过是一个相当寒酸的、看去全是一色普通村舍的小城镇，只有一个优点，就是或许所有的房子全是石结构的；但是墙上的灰泥早已掉光，砌墙的石块看来也开始剥落了。K蓦地想起了自己的家乡小镇；它同这座所谓的城堡相比几乎毫无逊色。如果K的目的只是观光，那么长途跋涉来到此地便太不值得，倒不如重访自己多年未归的故里更明智些呢。于是他在心中将家乡那座教堂的塔同眼前山上的塔作了一番比较。家乡那座塔从低到高、由粗而细笔直伸向天空，塔顶较宽，红瓦覆盖，诚然仍是一座不能超凡脱俗的建筑——我们凡人还能盖出什么别的样式来呢？——但它比大片低矮房屋有着更崇高的目标，比人们那晦暗的终日劳碌具有更明朗的意态。然而这里山上这座塔呢——这是此处唯一可见的塔，现在已能看清是一所住宅的塔，说不定就是城堡主楼的塔，它是一座单调的圆形建筑，部分蒙常春藤垂青加以覆盖，塔身有不少小窗子，此时在阳光下发出刺眼的反光——这使人觉得有些荒唐，塔顶类似阁楼，其雉堞瑟瑟缩缩、杂乱无章、残颓破败地戳向蓝天，就好像是一个害怕画错或是很不认真的孩子信手涂鸦胡乱画上去似的。面对这景象，人有一种感觉，仿佛这楼中有一个忧心忡忡的居民，按理本应老老实实地把自己关在楼中最偏僻的角落向隅而泣，现在居然冲破楼顶，探出身来向全世界亮相了。

K再次停下步来，似乎停下步子能增加他对眼前景象的判断能力。然而这时他受到了干扰。他站住的地方离村教堂（其实这不过是个小礼拜堂而已，只是像仓库似的扩建了一下，以便能容纳全村的教民）不远，教堂后面就是村子的小学。这学校是一座长长的低矮房子，奇怪的是它给人一种既是临时应急又非常古老的印象，楼前有一座围了铁栏杆、眼下是一片白雪覆盖着的花园。这时孩子们正好同老师一起从楼里走出来。一大群孩子把老师团团围住，每双眼睛都盯着他，从四面八方对他哇啦哇啦地唠叨个不住，K一点听不清他们那连珠炮似的话。这位老师是个年轻人，小个子，窄肩膀，不过样子并不可笑，身子直挺挺的，他从老远处一眼就看到了K，当然，除了他和他的一群学生之外，K是这方圆数里内目力所及唯一可见的人。K作为一个外乡人，首先打招呼问好，更何况是向这个在这里享有相当权威的小个子男人呢。“你好，老师！”他说。这话一出口，孩子们一下子就鸦雀无声，这突然的肃静作为他开口说话的先导，看样子使这位老师颇为得意。“您在观看城堡吗？”他问这问题时的态度，比K预料的要和蔼些，然而那语气似乎不赞同K这样做。“是的，”K说，“我不是本地人，是昨天晚上才到村里的。”——“您不喜欢这城堡吗？”教师很快又问。“什么？”K反问一句，有点震惊，然后用比较温和的语气重复那个问题：“我喜不喜欢这城堡吗？为什么您要猜想我不喜欢它呢？”——“没有一个外乡人喜欢它。”教师说。为了避免在这里说出一些不中听的话，K改变话题问：“您大概认识伯爵吧？”——“不认识。”教师说着便准备扭头走开。然而K并不罢休，他追问道：“什么？您不认识伯爵？”——“我怎么会认识他？”教师低声说，然后便用法语大声补充道：“请您考虑一下这里有这么些天真无知的孩子在旁边。”K听到这话便抓住时机问道：“老师，我可以拜访拜访您吗？我打算在这里待较长时间，可我现在就感到有点孤单了；我既不是农民中的一员，大概也不能算是城堡中的一员吧。”——“农民和城堡没有太大的区别。”教师说。“也许是这样，”K说，“可这丝毫改变不了我的处境。我可以拜访拜访您吗？”——“我住在天鹅巷肉铺。”虽然这话听起来只像是在说出住址，而不是发出邀请，K仍然说道：“好的，我一定来。”教师点点头，就同这时突然又大吵大嚷起来的孩子们一起继续走他的路了。不久，他们便消失在一条地势陡然下降的小巷里。

然而K现在却有些心神不宁，这次谈话使他心中颇感不快。自从到达此地以来，他现在第一次感到真正的疲倦了。到这里来的这一段很长的路，本来好像并不使他感觉吃力，这些天来，他是日复一日、从容不迫、一步一步地走过来的呵！——可是现在呢，过度紧张的后果到底还是显现出来了，无疑，来得真不是时候。他感到一种不可抗拒的强烈冲动，希望在这里尽量多结识一些人，但是每认识一个人都增添他一分倦意。〔1〕他想，在他今天这种竞技状态下，如果能咬牙坚持这次步行，至少走到城堡大门，那就已经是了不起的成绩了。

于是他又继续前行，但是路仍然很长。走着走着他发现，这条同时是村子主要街道的大路并不是通到城堡所在的山上去的，它只通到城堡近处，虽然眼看快到山脚下了，却像故意作弄人似的在那里拐了弯，然后，尽管沿它走下去并不会离城堡越来越远，却怎么也无法再接近它一步。K一直在期待着这条大路最终总会拐进城堡里去，也仅仅因为他抱有这一期望，他才不住地往前走；显然由于他感觉疲劳，才没有毅然决然离开它，另外，这个其长无比的村子也使他惊诧不已，它没有尽头，大路两边老是出现同样的小房子、冻了冰的窗户、厚厚的积雪，一个人影都不见——最后他终于还是下决心离开这条缠人磨人的大路，转进一条狭窄的小巷，这里雪更厚，把陷进雪里的脚拔出来艰巨异常，累得他满身大汗，突然间他站住了，再也走不动了。

细想起来，他倒也并不完全孤单，左右两旁不都是农舍吗？于是他捏了一个雪球，向一扇窗户掷去，门应声开了——这是他在村子里走了这大半天第一道打开的门，一个老农，身穿棕色皮夹克，头歪向一边，和蔼可亲地、颤巍巍地站在那里。“我可以到您屋里歇一会儿吗？”K说，“我太累了。”他完全没有听清老头说些什么，只是感激地接受了老者的好心：给他推过来一块木板，这木板当即把深陷雪中难以自拔的他救了出来，只跨过了三五步，他就进了老农屋里。
\par\vspace{0.3em}{\footnotesize\color{gray}（第一章）}
\newpage
\section{城堡来的少女}
\vspace{-0.3em}{\color{gray}\small\kaifont\noindent K 被老人从雪地里拉进了一间屋子。屋内光线昏暗，烟雾弥漫，有人洗澡，有人洗衣，有个女人靠在椅背上抱着婴儿。她说自己是“从城堡来的少女”。在这里，城堡不在山顶上，在每个人的行为方式里。}\vspace{0.5em}
\vspace{0.3em}

\noindent\rule{\textwidth}{0.4pt}

\vspace{0.3em}



屋子挺大，笼罩在一片朦胧的光线中。刚从外面进来的K一时什么也看不见。他趔趄几步，快撞到一个洗衣盆时，一只女人的手拉住了他。从屋子的一角传来好几个孩子的叫嚷声。从另一个角落则涌出滚滚烟雾，使本来就半明半暗的屋子变成一片漆黑。K如置身云雾中。“他喝醉了。”有人说。“您是谁？”另一个声音厉声喝问，然后大概是冲着老头问：“你为什么把他放进来？难道可以把大街小巷里东溜西窜的什么乱七八糟的人都放进来吗？”——“我是伯爵的土地测量员。”K说，试图用这话在那些直到现在仍唯闻其声不见其人的主人面前为自己辩护一下。“哦，原来是土地测量员。”这时响起的是一个女声，继之而来的是全然的寂静。“你们听说过我？”K问道。“当然。”同一个声音简短地答道。人家听说过他，这好像不是在抬举他吧。

浓烟终于稍稍散去了一些，K渐渐可以辨认四周了。看来今天是个大洗涤的日子。离门不远处有人在洗衣服，但雾气是来自屋子另一角，那里放着一个很大的木澡盆，K还从未见过这么大的澡盆——约有两张床那么大，两个男人泡在热气腾腾的水里洗澡。然而更加令人惊奇、可又弄不清究竟是什么东西使人惊奇的，是屋子右边那个角落。屋子后墙上有一个大缺口，通过这后墙上唯一的开口处，兴许是从院子里吧，射进来一道煞白的雪光，使一个坐在角落处一把高高的靠椅里、面带倦容的女人身上穿的衣服发出一种类乎丝绸的光亮。这女人怀里抱着一个婴儿，她四周有几个孩子在玩耍，显然都是农家孩子，但她似乎同他们并不属于一类人，当然，疾病和疲倦，也能使农民显得高雅些的。

“您请坐吧。”两个男人中的一个说，这人一脸络腮胡外加两髭向上翘起的小胡子，不断张着嘴喘气，他从水里伸出手，越过盆边——那副样子颇为滑稽——向一个大木箱指去，甩了K一脸热水。木箱上已经坐着一个人，正是那个把K放进来的老头，在那里发愣。对自己终于可以坐下来，K心里很是感激。现在没有人再理睬他了。在洗衣盆旁边的那个女人头发金黄，青春焕发，边洗边轻声哼着小曲，澡盆里的男人则不停地蹬腿翻身，孩子们想到他们跟前去，但一再被盆里溅出的大量洗澡水打退，K也未能幸免，靠椅里的女人泥塑木雕似的靠在椅背上，甚至连胸前的孩子她也不低头看上一眼，而只是视而不见地仰望空中。

K盯着她——这幅凝滞不动的、美丽、忧郁的雕像——大概看了好久，然后他一定是睡着了，因为当他听见有人大声叫他而猛然惊醒时，他的头是枕在他身旁的老者肩上的。这时两个男人早已洗完澡，穿好衣服站在K面前，澡盆里现在是孩子们在金发女人的看管下扑腾了。现在看来，两个男人中那个大嗓门的大胡子是地位低一些的。另一个的个头不比大胡子高，胡须也少得多，是个少言寡语、从容思考的人，他身材矮胖，长着一张宽脸，老是低着头。“土地测量员先生，”他说，“您不可以老待在这里，请原谅我的失礼。”——“我也不打算在这里待下去，”K说，“只是想稍稍休息一下。现在已经休息过了，我这就走。”——“您大概奇怪我们这里不太殷勤好客吧，”那人说，“我们这里可没有殷勤好客的习惯，我们是不需要客人的。”K小寐以后觉得精神稍好，注意力也比先前容易集中一些，听到这些直率的话很高兴。他行动不那么拘谨了，用手杖一会儿拄拄这，一会儿拄拄那，然后向靠椅里的女人走了过去。他也是这屋里个子最高的。

“不错，”K说，“你们要客人干什么呢。不过恐怕间或也会需要个把人的，比如我，土地测量员。”——“这个我不知道，”男人慢吞吞地说，“既然叫您来，那么大概是需要您，这可能是个例外，可我们呢，我们这些小老百姓，我们是按常规办事，这一点您不能怪我们。”——“哪里，哪里，”K说，“我对您只有感谢，对您和这里所有的人。”这时，出乎众人的意料，K突然嗖地一跳转过身去，站在那女人面前了。她用失神的蓝眼睛看着K，一块透明的真丝头巾一直盖到她前额正中，婴儿已在她怀里睡着了。“你是谁？”K问道。她拂袖——究竟这一轻蔑的表示是针对K的还是针对她自己的回答，弄不清楚——答道：“一个从城堡来的少女。”

这一切只是一瞬间的事。现在K发现自己已经一左一右夹在两个男人中间并被——好像根本没有别的交流思想的办法——一声不响地但却是用尽全力地拽到门边来。老头儿这时不知在乐什么，高兴得拍手。洗衣女人也在突然拼命吵嚷起来的孩子们旁边格格笑了。
\par\vspace{0.3em}{\footnotesize\color{gray}（第一章）}
\newpage
\section{电话与公文}
\vspace{-0.3em}{\color{gray}\small\kaifont\noindent K 来到村公所，村长卧病在床。他讲了两件事。打到城堡的电话会让最下一级部门的全部电话一起响起来，即便有人接听了也不能当真——那是“八音盒”。公文在各部门之间来回流转，每个环节都有回应，但最终没有裁断。如果说空间上的隔绝和人际关系里的渗透是你能感受到的痕迹，那制度上的错位呢？}\vspace{0.5em}
\vspace{0.3em}

\noindent\rule{\textwidth}{0.4pt}

\vspace{0.3em}



K对此地衙门的这种看法，首先在村长处得到了完全的证实。这是个和颜悦色、胡须刮得干干净净的胖子，他是在病中，在风湿病严重发作时躺在床上接见K的。“这么说您就是我们的土地测量员先生了。”他说，同时想坐起来表示欢迎，但怎么使劲也无用，于是带着歉意指指自己的腿又颓然躺回枕头上去。一个不苟言笑的女人给K拿来了一把椅子摆在床前，这间屋子的窗子本来就很小，又拉上了窗帘而更加昏暗，所以几乎只能看得见这女人那影影绰绰的轮廓。“您请坐，您请坐，土地测量员先生，”村长说，“请告诉我您有什么要求。”K念了克拉姆的信又加上几句自己的话。现在他再次体味到同衙门打交道时那种异常轻松的感觉。衙门简直就肩负着所有的重担，您可以把什么都推给它，自己落得个轻松自在。村长似乎从他那一方面同样感到了这一点，在床上很不舒服地翻了个身。最后他说：“您已经看到了，土地测量员先生，整个这件事情我是早就知道了的，我之所以还没有采取相应措施，第一是因为我生病，另外就是您很久没来，我已经在想您大概已经放弃这工作了。但现在呢，既然您亲临寒舍造访，我当然就必须把令人不快的全部实情倾囊相告。如您所说，您已被录用为土地测量员；但是遗憾得很，我们并不需要土地测量员。这里可以说一点让土地测量员做的工作也没有。我们这些小家小户的土地界线是划定了的，一应事项已全部有条不紊登记在案。产业易主的事极少发生，小的边界纠纷我们自己来解决。所以我们要土地测量员做什么？”虽然K原先完全没有细想过这一点，但他内心深处确信，自己是早就料到会得到类似说法了。〔7〕



正因为如此他立即说道：“我真万万没有想到事情是这样。这使我的全部打算都落空了。现在我只希望这是一个误会。”——“可惜没有误会，”村长说，“情况正像我说的那样。”——“可这也太不像话了！”K叫起来。“难道我千辛万苦长途跋涉到这里就为了让人再把我送回去吗！”——“这是另一码事，”村长说，“这件事我无权决定，但您所谓的那个误会是怎么产生的我倒可以给您解释解释。对一个像伯爵衙门这样的大衙门来说，有时难免会发生一个部门发出这一项指令、另一个部门发出另一项指令而互相不通气的事，虽然上一级的监督检查是极为严格的，但检查很自然地往往晚到一步，于是总会出点小差错，当然毫无例外都是在一些微乎其微的小事上出错，比如您这件事。在大事上我还没听说出过什么错，不过就是小事也够让人尴尬的了。至于说到您这件事，那么我可以坦率地把事情原原本本讲给您听，我用不着保守公务秘密——要保密我的官还太小，我是务农的，一辈子务农。很久以前，那时我当村长才几个月，上头来了一份公函，我记不清是哪个部门发出的了，那公函用高级长官特有的命令口气通知说，拟聘任一名土地测量员，指令村公所将有关该测量员工作必需的各项计划和全部文字材料立即备齐待命。这公函上指的人当然不可能是您，因为这是好多年前的事了，要不是现在我卧病在床，有的是时间去回想那些最最可笑的事情，那么我也是不会记起这事来的。”——“米齐，”他突然中断叙述对那个一直在屋里跑来跑去不知在忙乎什么的女人说，“请你去那边柜里看看，兴许能找到那份文件。”——“这项指令，”他接着对K解释，“是我刚当上村长那段时间下达的，那时我把所有的文件都保存起来了。”那女人当即打开柜子，K和村长在一旁看着她。柜子里各种文件纸张塞得满满当当的。柜门一开，两大捆像劈柴一样被捆成一大卷的文件便骨碌碌滚了出来，把那女人吓了一跳，一下子闪开了。“下面，可能在下面，”村长在床上指挥着，那女人则驯顺地张开双臂把大批大批的文件抱出来撂在地上，这样才好去拿最底下的那份文件。现在一卷卷文牍纸片已经堆满了半间屋子。“我们做了很多工作呵，”村长颔首说，“这些还只是一小部分呢。大宗的我都存放在库房里，而且，绝大部分文件都丢失了。谁有本事把全部文件都保存下来呀！库房里还有很多很多呢。”——“你找得着那份指令吗？”过一会儿他又对他女人说。“你先得找一份上面有‘土地测量员’字样，底下还画了蓝道道的文件。”——“这儿太暗了，”那女人说，“我去取一支蜡烛来，”说罢就迈步跨过地上的大堆文牍出屋去了。“我女人在这项繁重的公务中帮了我的大忙了，”村长说，“我还有别的重要事务，这项工作只能算是兼管一下。虽说我还有一个助手帮助我整理这些文字材料，这就是那位老师，可即使这样也仍然做不完，总是剩下许多堆积起来，还没整理的全部集中在那边那个柜子里，”他指指另一个柜子。“现在我这一病倒呢，那里简直就堆成山了，”说完这句他又疲倦地、然而却面带骄傲之色躺了下去。“我可不可以，”当那女人取来了蜡烛后又跪在柜子前继续找那份指令时，K说道，“帮着您太太一起找？”村长微笑着摇摇头：“刚才我说过，对您我没有什么公务秘密可保；可是，要让您本人去文件堆里找东西，那我确实又走得太远了。”此刻屋里一片安静，唯闻翻动纸张窸窸窣窣之声，在这种情况下村长甚至微微打起盹来。一阵轻轻的叩门声使K回身一看，当然又是两个助手。不过无论如何，两人现在总算有点长进：他们不是乒乒乓乓闯进屋来，而是站在外面通过微开的门缝轻声说：“我们在外面待着太冷了。”——“谁呀？”村长被惊醒了，问道。“没别人，是我的两个助手，”K说，“我不知道该让他们在哪儿等我，外面太冷，这儿嘛他们又让人讨厌。”——“我不在乎他们，”村长和气地说。“您让他们进来好了。其实我认识他们，我们是老相识了。”——“但是我讨厌他们，”K直言不讳，他把目光从两个助手移到村长身上，又从村长移回到助手身上，发现三个人的微笑一模一样，毫无区别。“既然你们两个到这儿来了，”他带着试探的口吻说，“那么就留下来帮着村长太太找一份文件吧，文件上写着‘土地测量员’，字底下画着蓝道道。”村长并无反对意见。就是说，不许K做的，两个助手倒可以做，两人也立即投身纸堆，甩开膀子大干起来，可是，与其说他们在找文件，倒不如说在纸堆里乱翻一气，而且，当一个人正读文件标题时，另一个人总是一把将文件从他手中抢走。那女人则跪在那现在已经空空如也的柜子前面，看样子根本没在继续找，至少现在蜡烛离她老远老远。



...



“您认识施瓦尔策吗？”K问。

“不认识，”村长说，“你大概认识吧，米齐？也不认识。不，我们不认识这个人。”

“这真是怪事，”K说，“他是城堡一位副主事的儿子。”

“亲爱的土地测量员先生，”村长说，“您怎么可以要求我认识所有城堡副主事的所有儿子？”

“很好，”K说，“那么您就只好相信我说的，承认他就是城堡一位副主事的儿子。同这个施瓦尔策，我在到这里的当天就有过一次不愉快的争论。后来他打电话给一个叫弗里茨的副主事询问，答复是我已经被聘任为土地测量员了。这您又怎么解释呢，村长先生？”

“很简单，”村长说，“这一点正好说明您从来还没有同我们的官府接触过嘛！所有这一切接触统统是表面文章，由于对情况完全无知，您竟将它们信以为真了。至于说到电话嘛，您瞧，我和官府在公务上的联系怎么说也是够多的了吧，可是我这里就没有电话。在酒店一类地方电话可能有相当大的好处，好比是八音盒那一类东西，然而更大的作用也就没有了。您在这里打过电话吗？打过？好，那您也许能明白我的意思。城堡里的电话显然非常灵；我听人说过，那里总是一个电话接着一个电话，这当然使工作的进度加快许多。那边连续不断的电话，我们这边电话机里听起来就是不停的嗡嗡声和歌唱声，这您一定也听到了。但我们这里的电话所能传达的唯一正确可信的东西，也就只是这种沙沙声和歌唱声罢了，别的信息全不是那么回事。我们同城堡之间并无确定的电话线路，没有总机接转我们打去的电话；从这里给城堡中某人挂电话时，那边最下一级的办事部门的电话就都一齐响起来，甚至于，这点我知道得很清楚，要不是城堡里几乎所有电话机上的响铃装置都关上了，那么所有的电话机就都会响起来的。偶尔有个别过于疲劳的官员觉着需要放松一下消遣消遣，特别是晚上或夜间，他就会把响铃装置打开；那时我们就能听到回话了，可那叫什么回话，不过是开开玩笑而已。这也是完全合情合理的。谁有权利为了想解决自己一点小小的私人困难而用电话铃声去干扰那些极为重要的、忙得人晕头转向的紧张工作？我也不明白为什么连一个外乡人也会以为，要是他比如说打电话给索尔蒂尼，那么接电话答复他的就一定是索尔蒂尼了。说多半会是另一个毫不相干的部门的一个小小记录员，这可能性不是反倒更大些吗？反过来说，在极少有的时候当然也会出现打电话给小记录员时反倒是索尔蒂尼本人来接电话的情况。那时当然是趁着还没听到他开腔说第一句话就赶紧跑开为妙。”
\par\vspace{0.3em}{\footnotesize\color{gray}（第五章）}
\newpage
\section{城市偏向（urban bias）：划一条区分“城”与“村”的线}
\vspace{-0.3em}{\color{gray}\small\kaifont\noindent 在发展中国家，城市与乡村之间那条线能不能划、怎么划，是理解一切城乡问题的前提。在城堡的世界里，我们面临一个类似的问题：城堡与村庄之间也同样找不到明确的边界。或许，我们不能因为没有边界就放弃分析。}\vspace{0.5em}
\vspace{0.3em}

\noindent\rule{\textwidth}{0.4pt}

\vspace{0.3em}



本章旨在界定“城市偏向”概念，并就其作为贫困与发展共存现象的罪魁祸首提出一个初步论证（后续章节将予以充实），同时附带说明它也是发展本身诸多低效与不足的成因。至此，我们已界定了“偏向”一词，并考察了两个我们期望据以证明大多数穷国存在“城市偏向”的规范：效率规范和公平规范。就前者而言，我们必须证明，如果将更大份额的资源投向农村，产出（按照第52页时间加权的含义）会提高。要依据公平规范来证明城市偏向，我们在上一节中已经论证过，大概只需证明农村人口比城市人口致富更慢（这在一定程度上是政策所致），以及农村人口起初就比城市人口贫穷很多，便已足够。后续章节将确立农村相对被剥夺的事实，考察农村资源以超出效率或公平合理限度的方式被转移至城市的情况，并检视导致这类转移的压力。



首先，我们必须仔细审视“城乡”二分法本身。可能有四种批评意见。（1）“这条线无法划定。城市在不知不觉中过渡为农村，中间经由城市份地、花园郊区与乡镇企业作为中介，界线模糊不清。”（2）“不止一条线。两个部门太少了。具体而言，存在一个第三类‘城市’部门，主要提供各类服务，在决定何种资源配置最优——或判断实际配置是否存在偏向时——必须加以考虑。”（3）“这条线应该画在别处：要画在首都与全国其他地方之间，而非画在城市与农村地区之间。”（4）“这条线画错了地图。资源配置与偏向存在于农业与工业之间，或消费品与投资品之间，或穷人与富人之间，而非农村与城市之间。”



**城乡界线**



在穷国，通常可以在城市与农村之间画出一条相当清晰的界线。少数例外不难想到：廉价且高度发达的交通系统使得兼职务农与进城通勤成为可能的地区，如斯里兰卡的湿润区；拥有可观小规模园艺产业的半发达第三世界城邦，如新加坡；市场与行政中心密布、近乎连续农村聚落的地区，如印度南部的喀拉拉邦。然而，一般而言，穷国既有城镇景观亦有乡村景观，二者之间的断裂是鲜明的。穷人很少有钱或有足够的热量能量进行长途通勤，住在农村（或住在城市）的人往往也就在当地工作。



第三世界农村地区的三种主要生活兼劳作模式是游牧、自营农场或庄园农场务农，以及定居村庄。在前两种情形中，从事露天农业的正是同一批人，他们也采用农村的生活方式。这使得农村部门的边界非常清晰。在第三种情形下，村庄与城镇之间的界线偶尔会模糊。一座城镇通常具有若干特征（紧凑性、人口密度、规模、非农业主导性）和诸多功能（市场、行政、社会、教育、交通），人口在三四千到七八千之间的某些地方满足部分标准但不满足另一些。不过，这个问题无论在什么意义上都是边缘性的。



穷国中绝大多数村庄社群的规模是数百人，主要是农业人口或为农业服务的工匠。大多数城镇人口居住在一万几千人以上的地方，极少从事农业或直接服务于农业的工作。因此，尽管各国人口普查对“城市”的定义看似杂乱多样，但对于绝大多数地方和居民而言，城乡二分法依然是清晰且离散的。本书所诊断的偏向源于人口一万到两万及以上规模的较大城镇，并使这类城镇从中受益。



\vspace{0.3em}

\noindent\rule{\textwidth}{0.4pt}

\vspace{0.3em}
\par\vspace{0.3em}{\footnotesize\color{gray}（Lipton, M. (1977). \textit{Why poor people stay poor: Urban bias in world development}. Temple Smith. 由 DeepSeek V4 Pro 翻译。）}
\newpage
\begin{center}
{\LARGE \bfseries Day 4\\[1em]共谋、非正式制度与组织行为}
\end{center}
\addcontentsline{toc}{section}{\protect\numberline{17}共谋、非正式制度与组织行为}
\stepcounter{section}

\vspace{3em}
\vspace{-0.3em}{\color{gray}\small\kaifont \noindent 理解组织行为有两种常见路径。一是看正式制度，谁管谁，什么事走什么流程。二是看个人素质，这个人是不是偷懒了，那个人是不是搞腐败了。两种路径都有解释力，但都漏掉了一类现象。有些行为既不是规则上写着的，也不归因于哪个人的错误。它稳定存在，反复再生，组织中的所有人都知道它在发生，也都默认它是正常的。周雪光把这类现象称为“制度化了的非正式行为”。他以中国基层政府间“上有政策、下有对策”的共谋现象为切入口，论证了一个在组织研究中具有普遍意义的命题。正式制度本身会催生出大量非正式行为，这些行为不是制度的例外，而是制度的实现形式。}\vspace{0.5em}
\vspace{0.3em}

\noindent\rule{\textwidth}{0.4pt}

\vspace{0.3em}



[内容提要]中国政府行为的一个突出现象是，在执行来自上级部门特别是中央政府的各种指令政策时，基层上下级政府常常共谋策划、暗渡陈仓，采取“上有政策、下有对策”的各种手段，来应付这些政策要求以及随之而来的各种检查，导致了实际执行过程偏离政策初衷的结果。本文从组织学角度，对这类现象提出一个理论解释。本文的中心命题是：在中国行政体制中，基层政府间的共谋行为已经成为一个制度化了的非正式行为；这种共谋行为是其所处制度环境的产物，有着广泛深厚的合法性基础。本文讨论政府组织制度的三个悖论，对这一现象提出理论解释：一是政策一统性与执行灵活性的悖论，二是激励强度与目标替代的悖论，三是科层制度非人格化与行政关系人缘化的悖论。本研究强调，共谋行为不能简单地归咎于政府官员或执行人员的素质或能力，其稳定存在和重复再生是政府组织结构和制度环境的产物；是现行组织制度中决策过程与执行过程分离所导致的结果；在很大程度上也是近年来政府制度设计特别是集权决策过程和激励机制强化所导致的非预期结果。而欲改变这一状况，首先需要对政府组织现象进行深入系统的研究，提出有力的理论解释。



一、现象与问题



作为中央政策执行过程的最终环节，基层政府在中国政府组织制度中有着极为重要的战略位置。已有许多研究工作关注分析基层政府的状况和行为（参见赵树凯，2005，2006b；在社会学领域，参见周飞舟，2006；张静，2000；吴毅，2007）。在中国政府组织制度中，贯彻执行上级政府的各项政策指令是基层政府的一个重要工作内容。上级决策者和政策研究者经常用“上传下达、令行禁止”这样的词汇来形容基层政府工作的理想状态。但实际生活中一个引人注目的现象是，在执行来自上级部门特别是中央政府的各种政策指令时，一些基层政府常常共谋策划，采取“上有政策、下有对策”的各种做法，联手应付自上而下的政策要求以及随之而来的各种检查，导致了实际执行过程偏离政策初衷的结果。例如，一位镇政府主管计生工作的官员这样描述有关的检查工作：



\begin{quotation}\kaifont\noindent 省市县都组织检查。当省里来检查时，市、县、镇联合起来对付省的检查团；当市里来检查时，县、镇就联合起来对付市里的检查团。省里到县里检查时，事先不通知去哪个镇或村。但是，实际上各地政府都联合起来对付检查团。省里来的时候，市里、县里就提前来打招呼：“你们那里的问题处理好……”检查团到了县里，各个镇都得到通知，严阵以待。一旦领队的得到通知去哪个村，在去那个村的路上就有电话通知到那个村，连车号、时间、地点的消息都十分详尽。检查团一般都是早晨8点以前就到达。所以，一大早，村里妇女小组长就在各个路口把守，一旦发现情况就立即通知村里，把有问题的孩子转移出去……我工作6年来，除了有一个乡被省检查团查出来外，只有两次（两个镇）是县里检查出计划外的问题，做了一票否决的处理，镇书记和镇长都调换工作，没有重用。县里也没有向上报告。（访谈记录 H2346）\end{quotation}



这种上下级政府间联手应对更上级政府的现象普遍存在。媒介报道更多地注意到了这些负面状况。例如，煤窑安全事故发生后，事发地基层政府间层层包庇，封锁道路，某些基层官员滥用权力，压迫民众；一旦事发，自上而下调查时，那些基层政府上下级也经常互相掩盖，大事化小，小事化了（参见陈桂棣、春桃，2004；高王凌，2006）。这些行为有着各种表现形式，或者是基层政府出现问题而游说直接上级部门为之掩盖保护；或者是直接上级部门要求下级部门掩盖问题，以应付来自更上级政府的要求和检查。更多的情形下，这类行为体现在上下级基层政府持续双向的互动过程中。因为这些应对策略和行为常常与上级政策指令不符甚至相悖，所以它们大多是通过非正式方式加以实施（例如，“打招呼”、私下安排等）。本文用基层政府间“共谋行为”这一概念来概括描述这类非正式行为。



值得强调的是，本文所关注的基层政府“共谋行为”是在中性意义上使用的一个分析概念。如下文所讨论的，这些“共谋现象”在很多情况下是基层政府执行政策时所表现的灵活性和适应性，是制度创新的一个重要源泉。而那些负面意义上的基层政府间共谋行为，虽然是中央政府所不愿意看到并且力图克服的，但这类行为的范围之广、程度之深、表现之公开、运作之坚韧却是少见的。在中国社会转型过程中，随着政府制度职能的演变和改革，这一问题也随之日益突出，甚至有愈演愈烈的趋势。重复再现的组织现象是建筑在稳定持续的组织制度基础之上和相应的组织环境之中的。本文对这一现象从组织学的角度尝试提出一个理论解释。



我们的中心命题是：在中国行政体制中，基层政府间的共谋行为已经成为一个制度化了的非正式行为。即，这种共谋行为是基层政府所处制度环境的产物，因此有着广泛深厚的合法性基础和特定的制度逻辑。这里使用的“制度化”概念来自组织社会学中的新制度学派（Meyer \& Rowan，1977；参见周雪光，2003）。“制度化”的意义是，某个特定的组织行为或组织结构被社会或组织环境中有关各方所广泛接受，因此具有合法性或合理性；而组织环境中的制度化机制促就这类现象的稳定存在、重复再生。从这个角度来看，基层政府间的共谋现象恰具这些制度化特点：一方面，这些共谋行为大多与正式政策法令的明文规定不一致甚至相悖，旨在应付、欺骗、蒙蔽上级政府，因此多以非正式形式出现。但另一方面，这些行为并不是个别官员或个别部门的私人活动；在许多情形下，它们在正式组织权力结构下公开运作，以政府部门的组织权威辅以实施，甚至是大张旗鼓地加以部署安排。而这些做法已经成为上下级政府甚至中央政府的“共有常识”（common knowledge）。正是在这个意义上，我们说基层政府间共谋现象是特定的组织环境中制度化了的非正式行为。
\par\vspace{0.3em}{\footnotesize\color{gray}（周雪光. (2009). 基层政府间的“共谋现象”——一个政府行为的制度逻辑. 开放时代, (12), 40-55.）}
\newpage
\begin{center}
{\LARGE \bfseries Day 5\\[1em]教育作为一种行政现象}
\end{center}
\vspace{3em}
\section{高考改卷组织与确定性偏好}
\vspace{-0.3em}{\color{gray}\small\kaifont \noindent 你们即将读到两种材料。一种是制度文件。《国家教育统一考试网上评卷工作考务管理办法》，教育部办公厅发布的。它规定了高考改卷的组织架构、人员选聘资格、评分细则制定流程、质量控制机制。它是一套精密机器的技术说明书。另一种是阅卷教师写的经验分享。它告诉你在这套机器里面工作是什么感觉。

\noindent 制度设计者说“这把尺子量到底是为了公平”，身处机器中的人说“我根本没时间欣赏你的文采”。两段话拼在一起，你看到的不再是公平或残酷的单一结论，而是一个组织行为问题：当“减少误差”成为唯一的优化目标，这套机制鼓励的不是最优答案，而是最可预测的答案。}\vspace{0.5em}
\vspace{0.3em}

\noindent\rule{\textwidth}{0.4pt}

\vspace{0.3em}

{\color{gray}\small （教育部办公厅，2006，2006年国家教育统一考试网上评卷工作考务管理办法）}



第三章　评卷工作的组织管理



第七条　评卷工作由省级教育考试机构负责组织和管理，有关高等学校给予支持配合。



第八条　评卷工作由省级教育考试机构在指定的高等学校或指定的其它场所集中进行。



第九条　评卷场地应有专人保卫，评卷员和评卷工作人员出入须凭《评卷工作证》，不得随意变更，也不得将《评卷工作证》转借他人。



评卷员使用的密码只限本人掌握，不得告知其他评卷员。



第十条　应成立由省级教育考试机构、评卷高校和学科负责人参加的评卷领导小组。



各学科（课程）成立评卷小组，负责人由本学科（课程）具有副教授以上职称的教师担任。



评卷程序和复评数量、标准等由评卷领导小组确定。



第十一条　评卷员由评卷领导小组聘任。评卷员应是从事本学科（课程）教学工作、业务水平较高、责任心强、作风正派的教师。评阅作文、论述等主观性试题的教师，必须具有本学科中、高级以上职称。



第十二条　评卷前应对全体评卷员及有关工作人员进行保密纪律的教育和评卷规定、评卷程序的培训。如果省级教育考试机构直接在社会上聘用评卷员，应与所聘任的评卷员签订评卷协议。



第十三条　评卷前各学科（课程）评卷小组负责人应组织本学科（课程）评卷员认真研究试题、答案和评分参考，并进行试评。在试评基础上制订评分细则，评分细则应报省级教育考试机构备案。



第十四条　评卷员必须严格执行评分细则，不得擅自更改。凡在评卷过程中对评分细则有异议，属于全国统一命题的科目的，须由省级教育考试机构报教育部考试中心（自考办）同意后方可改动；属于省自主命题的科目的，须报省级教育考试机构同意后方可改动。



第十五条　学科评卷小组负责人应熟悉评卷系统。评卷过程中，学科评卷小组负责人应随时掌握评卷进度和评卷员的评卷质量。对于不能准确把握评分细则的评卷员，经提醒后仍不能纠正的应及时进行调整。



第十六条　省级教育考试机构应与为网上评卷提供技术支持和服务的有关公司及所聘任的评卷工作人员签订评卷工作协议。









{\color{gray}\small （语文帮高中版，2026，高考阅卷不看文采、不看感悟，只看踩点给分（新课标一卷通用）}，今日头条）



一、先懂底层逻辑：高考阅卷，是标准化流水线，不是文学鉴赏





高考是国家级选拔性标准化考试，第一原则是：绝对公平、统一标准、可量化打分。





无数学生搞错了阅卷本质：



你以为阅卷老师是语文名师，慢慢品读你的答案、共情你的感悟、欣赏你的文笔；



真实阅卷现状： 每位阅卷老师日均阅卷3000–5000份，单道主观题阅卷时长3–8秒。





在极致的阅卷速度、统一的评分规则下：





1. 优美的辞藻、流畅的段落、细腻的个人感悟，无法量化，一律不计分；



2. 阅卷老师没有时间品读情感、赏析文笔，只做一件事：扫描关键词、核对得分点；



3. 所有主观题分数，全部绑定固定答题角度、专属得分术语、文本对应依据，缺一不得分。





简单一句话：高考语文的分数，是按点结算的，不是按情分、文笔分的。





你的个人感悟再深刻、文字再优美，只要没答出官方规定的得分点，就是0分；



你的答案朴实直白、毫无文采，只要精准踩中所有关键词、答题维度，就是满分。



{\color{gray}\small （天意，2026，高考阅卷老师的压力总结，小红书）}



额外的淘汰机制，全程都不能松懈

上岗前必须通过试评考核，打分和专家组标准误差超标，根本没有正式阅卷资格；阅卷中途一旦稳定性不达标，会被中途清退，这也让老师全程不敢放松心态。

简单总结：阅卷不是简单改卷子，相当于在严密的监控下，连续一周多做高精密的精细活，既要快、又要准，还要严守纪律，心理和身体都会达到负荷上限。



{\color{gray}\small （冰川思想库，2019，那些年，我经历过的高考阅卷，知乎）}



高考阅卷更算不得神圣。为什么参加高考阅卷？大多数人就是为了挣钱。





高考阅卷一般要连续三到四天，每天从早到晚，是比较辛苦的。老教师一般不愿参加，就是觉得累，挣钱还不多。那么，能挣多少钱呢？当时（十几年前）一次高考阅卷大概能拿到一千多元。





因为缺人，每次都有一定数量的研究生参与进来，对一些经济条件不好的研究生来说，几天时间能挣个千儿八百的还不错，据说，有的研究生还会特意争取阅卷的机会。





当然，没有高考阅卷经验的研究生的阅卷质量往往被重点关注，反复核查。具体到某个人高考阅卷的收入，主要是由产量（阅卷量）决定的。





每个人阅卷的数量都由专人统计。眼疾手快的人，可就占便宜了，阅卷收入要远远高于动作迟缓的人。

\newpage\section{街头官僚的职责拮抗}
\vspace{-0.3em}{\color{gray}\small\kaifont \noindent 前面我们看了高考改卷组织如何令可能怀有良好意愿的教师化为追求确定性的给分机器。你可能会问：那老师教课的时候呢？老师关心学生的时候呢？那些不在改卷体系管辖范围内的日常互动，难道也被“标准化”了吗？

\noindent 答案是没有。而这恰恰是问题所在。

\noindent 李普斯基在上世纪八十年代提出了“街头官僚”这个概念。他说的不是街上的公务员。他说的是那些在日常工作中直接面对公民、并对他人生活产生实际影响的一线工作者——警察、社工、法官、教师。这些人的核心特征不是职位高低，而是他们在执行政策的过程中拥有“不可消除的自由裁量权”——因为政策不可能覆盖人与人的每一次具体互动。

\noindent 这种自由裁量权带来的后果是什么？Lipsky 说，街头官僚同时承担着中心职责和边缘职责。中心职责是组织给他的任务，被上级严密管理、频繁考核。边缘职责则分散在日常互动的间隙中——它们不容易被识别、量化和管理，但同样真实地影响着他人的生活质量。这两类职责之间存在拮抗关系：管好了一个，另一个可能就会被压掉。}\vspace{0.5em}
\vspace{0.3em}

\noindent\rule{\textwidth}{0.4pt}

\vspace{0.3em}





街头官僚工作环境的一个根本特征是：他们必须面对服务对象对自己决定的个人反应，无论他们如何处理这些决定带来的后果。说人的自我评价受到街头官僚行为的影响，就等于说人对政策是反应性的。这不仅限于潜意识过程。街头官僚机构的服务对象会对真实的或感知到的不公正感到愤怒，会制定策略来讨好工作人员，会因街头官僚的决定而表现出感激和兴高采烈，或者闷闷不乐和消极被动。被电话公司、车管局或其他政府机构以一种冷漠和常规的方式对待是一回事——这些机构的代理人对你的具体情况一无所知。而被一个你正在与之面对面交谈、并期望至少得到一次开放而富有同情心的倾听的人以“官僚式的”（在这个词的贬义意义上）方式进行筛选、分类和对待，则完全是另一回事。简而言之，街头官僚的工作现实与官僚制所崇尚的非人格化超脱决策理想相去甚远。恰恰相反，在街头官僚机构中，关键决定的对象——人——实际上会因为这些决定而发生改变。



街头官僚也是公民反应的焦点，因为他们的自由裁量权开启了一种可能性：他们可能替某人做出有利的回应。他们对“公共利益”所承担的宽泛而模糊的义务，让人们怀抱希望，期待这个具体的工作人员会对自己的处境抱有好意或偏向。因此，在一个由庞大而非人格化的机构组成的世界中——这些机构表面上掌握着通往重要福利、制裁和机会的钥匙——工作定义的模糊性维系着“在这里或许能找到一个肯帮忙的人”的希望。
\par\vspace{0.3em}{\footnotesize\color{gray}（Lipsky, M. (1980). Street-Level Bureaucracy: Dilemmas of the Individual in Public Services. Russell Sage Foundation. 由 DeepSeek V4 Pro 翻译。）}
\newpage\section{教育作为排他性的国家治理领域}



\vspace{-0.3em}{\color{gray}\small\kaifont \noindent 1985年5月27日，中共中央发布了《关于教育体制改革的决定》。这是一个里程碑式的纲领性文件。它开篇就说：“教育必须为社会主义建设服务，社会主义建设必须依靠教育。”它把教育定位为经济建设的人才供应线，要求培养“数以亿计”的劳动者、“数以千万计”的技术人员和各方面的工作者。教育行政组织与生活在学生中的教师的复杂心绪，从那一刻起，就被嵌入了一个更大的国家目标：快速实现现代化的经济和行政建制。}\vspace{0.5em}



\vspace{0.3em}

\noindent\rule{\textwidth}{0.4pt}

\vspace{0.3em}





{\color{gray}\small （中共中央，1985，中共中央关于教育体制改革的决定）}



党的十二届三中全会关于经济体制改革的决定，为我国社会生产力的大发展、为我国社会主义物质文明和精神文明的大提高，开辟了广阔的道路。今后事情成败的一个重要关键在于人才，而要解决人才问题，就必须使教育事业在经济发展的基础上有一个大的发展。

教育必须为社会主义建设服务，社会主义建设必须依靠教育。社会主义现代化建设的宏伟任务，要求我们不但必须放手使用和努力提高现有的人才，而且必须极大地提高全党对教育工作的认识，面向现代化、面向世界、面向未来，为九十年代以至下世纪初叶我国经济和社会的发展，大规模地准备新的能够坚持社会主义方向的各级各类合格人才。要造就数以亿计的工业、农业、商业等各行各业有文化、懂技术、业务熟练的劳动者。要造就数以千万计的具有现代科学技术和经营管理知识，具有开拓能力的厂长、经理、工程师、农艺师、经济师、会计师、统计师和其他经济、技术工作人员。还要造就数以千万计的能够适应现代科学文化发展和新技术革命要求的教育工作者、科学工作者、医务工作者、理论工作者、文化工作者、新闻和编辑出版工作者、法律工作者、外事工作者、军事工作者和各方面党政工作者。所有这些人才，都应该有理想、有道德、有文化、有纪律，热爱社会主义祖国和社会主义事业，具有为国家富强和人民富裕而艰苦奋斗的献身精神，都应该不断追求新知，具有实事求是、独立思考、勇于创造的科学精神。这就向我国教育事业的发展和教育体制的改革，提出了伟大而又艰巨的任务。

建国以来，我国教育事业的发展走过了曲折的道路。经过解放初期的接管改造和以高等学校院系调整为中心的教育改革，我们把旧中国的半殖民地半封建教育事业转变成为社会主义教育事业。三十几年来，依靠广大教育工作者的辛勤努力，教育事业取得了中国历史上从来没有过的巨大的发展，成绩是显著的。今天战斗在我国各条战线上的广大有文化的劳动者和各方面工作的骨干力量，绝大部分都是建国以后培养出来的。但是，另一方面，从五十年代后期开始，由于全党工作重点一直没有转移到经济建设上来，由于“阶级斗争为纲”的“左”的思想的影响，教育事业不但长期没有放到应有的重要地位，而且受到“左”的政治运动的频繁冲击。“文化大革命”更使这种“左”的错误走到否定知识、取消教育的极端，从而使教育事业遭到严重破坏，广大教育工作者遭受严重摧残，耽误了整整一代青少年的成长，并且使我国教育事业同世界发达国家之间在许多方面本来已经缩小的差距又拉大起来。
\newpage
\begin{center}
{\LARGE \bfseries Day 6\\[1em]组织的纹理}
\end{center}
\vspace{3em}
\section{前哨的第六封信}
\vspace{-0.3em}{\color{gray}\small\kaifont \noindent K 在村子里待了足够久。白天和村长谈话，晚上在小屋里整理自己的所见。他不再试图从外部观察城堡，也不再追问“它到底在哪”。他把这几天看到的几件事拼在一起，写给即将出发的队友。}\vspace{0.5em}
\vspace{0.3em}

\noindent\rule{\textwidth}{0.4pt}

\vspace{0.3em}



同志们：



你们听完了周雪光，也讨论过 Lipsky——现在我想带你们回到城堡本身。



我在村里看到了三样东西。公文在每个环节都有答复，但从没有裁断：下级负责遗忘，上级负责铭记。职能之间是彻底区隔的，一个部门不可以问另一个部门“是不是你发的”——这是纹理，是组织在生长和断裂中不变的东西。然后巴纳巴斯去了办公厅，穿着自己的衣服，站在那些他能越过的和不能越过的栅栏之间，发现它们看上去一模一样。



这三件事是同一件事。你们先读，读完请告诉我你们看到了什么。



K

\newpage
\section{电话与公文}
\vspace{-0.3em}{\color{gray}\small\kaifont \noindent K 来到村公所，村长卧病在床。他讲了两件事。打到城堡的电话会让最下一级部门的全部电话一起响起来，即便有人接听了也不能当真——那是“八音盒”。公文在各部门之间来回流转，每个环节都有回应，但最终没有裁断。如果说空间上的隔绝和人际关系里的渗透是你能感受到的痕迹，那制度上的错位呢？}\vspace{0.5em}
\vspace{0.3em}

\noindent\rule{\textwidth}{0.4pt}

\vspace{0.3em}



K对此地衙门的这种看法，首先在村长处得到了完全的证实。这是个和颜悦色、胡须刮得干干净净的胖子，他是在病中，在风湿病严重发作时躺在床上接见K的。“这么说您就是我们的土地测量员先生了。”他说，同时想坐起来表示欢迎，但怎么使劲也无用，于是带着歉意指指自己的腿又颓然躺回枕头上去。一个不苟言笑的女人给K拿来了一把椅子摆在床前，这间屋子的窗子本来就很小，又拉上了窗帘而更加昏暗，所以几乎只能看得见这女人那影影绰绰的轮廓。“您请坐，您请坐，土地测量员先生，”村长说，“请告诉我您有什么要求。”K念了克拉姆的信又加上几句自己的话。现在他再次体味到同衙门打交道时那种异常轻松的感觉。衙门简直就肩负着所有的重担，您可以把什么都推给它，自己落得个轻松自在。村长似乎从他那一方面同样感到了这一点，在床上很不舒服地翻了个身。最后他说：“您已经看到了，土地测量员先生，整个这件事情我是早就知道了的，我之所以还没有采取相应措施，第一是因为我生病，另外就是您很久没来，我已经在想您大概已经放弃这工作了。但现在呢，既然您亲临寒舍造访，我当然就必须把令人不快的全部实情倾囊相告。如您所说，您已被录用为土地测量员；但是遗憾得很，我们并不需要土地测量员。这里可以说一点让土地测量员做的工作也没有。我们这些小家小户的土地界线是划定了的，一应事项已全部有条不紊登记在案。产业易主的事极少发生，小的边界纠纷我们自己来解决。所以我们要土地测量员做什么？”虽然K原先完全没有细想过这一点，但他内心深处确信，自己是早就料到会得到类似说法了。〔7〕



正因为如此他立即说道：“我真万万没有想到事情是这样。这使我的全部打算都落空了。现在我只希望这是一个误会。”——“可惜没有误会，”村长说，“情况正像我说的那样。”——“可这也太不像话了！”K叫起来。“难道我千辛万苦长途跋涉到这里就为了让人再把我送回去吗！”——“这是另一码事，”村长说，“这件事我无权决定，但您所谓的那个误会是怎么产生的我倒可以给您解释解释。对一个像伯爵衙门这样的大衙门来说，有时难免会发生一个部门发出这一项指令、另一个部门发出另一项指令而互相不通气的事，虽然上一级的监督检查是极为严格的，但检查很自然地往往晚到一步，于是总会出点小差错，当然毫无例外都是在一些微乎其微的小事上出错，比如您这件事。在大事上我还没听说出过什么错，不过就是小事也够让人尴尬的了。至于说到您这件事，那么我可以坦率地把事情原原本本讲给您听，我用不着保守公务秘密——要保密我的官还太小，我是务农的，一辈子务农。很久以前，那时我当村长才几个月，上头来了一份公函，我记不清是哪个部门发出的了，那公函用高级长官特有的命令口气通知说，拟聘任一名土地测量员，指令村公所将有关该测量员工作必需的各项计划和全部文字材料立即备齐待命。这公函上指的人当然不可能是您，因为这是好多年前的事了，要不是现在我卧病在床，有的是时间去回想那些最最可笑的事情，那么我也是不会记起这事来的。”——“米齐，”他突然中断叙述对那个一直在屋里跑来跑去不知在忙乎什么的女人说，“请你去那边柜里看看，兴许能找到那份文件。”——“这项指令，”他接着对K解释，“是我刚当上村长那段时间下达的，那时我把所有的文件都保存起来了。”那女人当即打开柜子，K和村长在一旁看着她。柜子里各种文件纸张塞得满满当当的。柜门一开，两大捆像劈柴一样被捆成一大卷的文件便骨碌碌滚了出来，把那女人吓了一跳，一下子闪开了。“下面，可能在下面，”村长在床上指挥着，那女人则驯顺地张开双臂把大批大批的文件抱出来撂在地上，这样才好去拿最底下的那份文件。现在一卷卷文牍纸片已经堆满了半间屋子。“我们做了很多工作呵，”村长颔首说，“这些还只是一小部分呢。大宗的我都存放在库房里，而且，绝大部分文件都丢失了。谁有本事把全部文件都保存下来呀！库房里还有很多很多呢。”——“你找得着那份指令吗？”过一会儿他又对他女人说。“你先得找一份上面有‘土地测量员’字样，底下还画了蓝道道的文件。”——“这儿太暗了，”那女人说，“我去取一支蜡烛来，”说罢就迈步跨过地上的大堆文牍出屋去了。“我女人在这项繁重的公务中帮了我的大忙了，”村长说，“我还有别的重要事务，这项工作只能算是兼管一下。虽说我还有一个助手帮助我整理这些文字材料，这就是那位老师，可即使这样也仍然做不完，总是剩下许多堆积起来，还没整理的全部集中在那边那个柜子里，”他指指另一个柜子。“现在我这一病倒呢，那里简直就堆成山了，”说完这句他又疲倦地、然而却面带骄傲之色躺了下去。“我可不可以，”当那女人取来了蜡烛后又跪在柜子前继续找那份指令时，K说道，“帮着您太太一起找？”村长微笑着摇摇头：“刚才我说过，对您我没有什么公务秘密可保；可是，要让您本人去文件堆里找东西，那我确实又走得太远了。”此刻屋里一片安静，唯闻翻动纸张窸窸窣窣之声，在这种情况下村长甚至微微打起盹来。一阵轻轻的叩门声使K回身一看，当然又是两个助手。不过无论如何，两人现在总算有点长进：他们不是乒乒乓乓闯进屋来，而是站在外面通过微开的门缝轻声说：“我们在外面待着太冷了。”——“谁呀？”村长被惊醒了，问道。“没别人，是我的两个助手，”K说，“我不知道该让他们在哪儿等我，外面太冷，这儿嘛他们又让人讨厌。”——“我不在乎他们，”村长和气地说。“您让他们进来好了。其实我认识他们，我们是老相识了。”——“但是我讨厌他们，”K直言不讳，他把目光从两个助手移到村长身上，又从村长移回到助手身上，发现三个人的微笑一模一样，毫无区别。“既然你们两个到这儿来了，”他带着试探的口吻说，“那么就留下来帮着村长太太找一份文件吧，文件上写着‘土地测量员’，字底下画着蓝道道。”村长并无反对意见。就是说，不许K做的，两个助手倒可以做，两人也立即投身纸堆，甩开膀子大干起来，可是，与其说他们在找文件，倒不如说在纸堆里乱翻一气，而且，当一个人正读文件标题时，另一个人总是一把将文件从他手中抢走。那女人则跪在那现在已经空空如也的柜子前面，看样子根本没在继续找，至少现在蜡烛离她老远老远。



...



“您认识施瓦尔策吗？”K问。

“不认识，”村长说，“你大概认识吧，米齐？也不认识。不，我们不认识这个人。”

“这真是怪事，”K说，“他是城堡一位副主事的儿子。”

“亲爱的土地测量员先生，”村长说，“您怎么可以要求我认识所有城堡副主事的所有儿子？”

“很好，”K说，“那么您就只好相信我说的，承认他就是城堡一位副主事的儿子。同这个施瓦尔策，我在到这里的当天就有过一次不愉快的争论。后来他打电话给一个叫弗里茨的副主事询问，答复是我已经被聘任为土地测量员了。这您又怎么解释呢，村长先生？”

“很简单，”村长说，“这一点正好说明您从来还没有同我们的官府接触过嘛！所有这一切接触统统是表面文章，由于对情况完全无知，您竟将它们信以为真了。至于说到电话嘛，您瞧，我和官府在公务上的联系怎么说也是够多的了吧，可是我这里就没有电话。在酒店一类地方电话可能有相当大的好处，好比是八音盒那一类东西，然而更大的作用也就没有了。您在这里打过电话吗？打过？好，那您也许能明白我的意思。城堡里的电话显然非常灵；我听人说过，那里总是一个电话接着一个电话，这当然使工作的进度加快许多。那边连续不断的电话，我们这边电话机里听起来就是不停的嗡嗡声和歌唱声，这您一定也听到了。但我们这里的电话所能传达的唯一正确可信的东西，也就只是这种沙沙声和歌唱声罢了，别的信息全不是那么回事。我们同城堡之间并无确定的电话线路，没有总机接转我们打去的电话；从这里给城堡中某人挂电话时，那边最下一级的办事部门的电话就都一齐响起来，甚至于，这点我知道得很清楚，要不是城堡里几乎所有电话机上的响铃装置都关上了，那么所有的电话机就都会响起来的。偶尔有个别过于疲劳的官员觉着需要放松一下消遣消遣，特别是晚上或夜间，他就会把响铃装置打开；那时我们就能听到回话了，可那叫什么回话，不过是开开玩笑而已。这也是完全合情合理的。谁有权利为了想解决自己一点小小的私人困难而用电话铃声去干扰那些极为重要的、忙得人晕头转向的紧张工作？我也不明白为什么连一个外乡人也会以为，要是他比如说打电话给索尔蒂尼，那么接电话答复他的就一定是索尔蒂尼了。说多半会是另一个毫不相干的部门的一个小小记录员，这可能性不是反倒更大些吗？反过来说，在极少有的时候当然也会出现打电话给小记录员时反倒是索尔蒂尼本人来接电话的情况。那时当然是趁着还没听到他开腔说第一句话就赶紧跑开为妙。”
\par\vspace{0.3em}{\footnotesize\color{gray}（第五章）}
\newpage
\section{土地测量员聘用始末}
\vspace{-0.3em}{\color{gray}\small\kaifont \noindent K 不只要一个简单的说法。他要搞清楚，聘他来做土地测量员的这个决定到底是怎么跑出来的。村长讲了很久。故事的核心不是答案。}\vspace{0.5em}
\vspace{0.3em}

\noindent\rule{\textwidth}{0.4pt}

\vspace{0.3em}



“这么说这两个助手，”村长又开口了，洋洋得意地微笑着，似乎在说：所有这些全是他一手安排的，而谁都无法预想到这一点，“这么说，他们让您感到讨厌，可他们是您自己的助手呵。”——“不对，”K冷冷地说，“他们是我到此地后才跑到我身边来的。”——“怎么，跑来的，”村长说，“您的意思大概是说派来的吧。”——“好吧，就算是派来的吧，”K说。“但说他们是从天外飞来的也未尝不可，这种派法也够欠考虑的了。”——“这里办什么事都不会欠考虑，”村长说，急得连腿疼也忘了，一骨碌坐起来。“办什么事都不欠考虑，”K说，“那么招聘我这事又怎么说？”——“聘用您也是经过深思熟虑的，”村长说，“只是有一些次要因素起了干扰作用罢了，我可以用有关的文件向您证明这一点。”——“这些文件恐怕找不到了，”K说。“找不到？”村长叫起来。“米齐，请快点找！不过，没有文件我也可以先把事情的经过讲给您听听。刚才我提到的那份指令，我们接到后当即回复上头表示感谢，说我们不需要土地测量员。但是这封回信看来是没有送回发出指令的那个部去，我管它叫A部吧，而是误送到另一个部，送到B部去了。这样A部就一直没有得到回音，然而可惜连B部也没有收到我们复信的全部；一种可能是公文袋里装的信落在村公所，再一种可能是在半路上遗失了——肯定不会在部里遗失，这我敢担保，总之，B部也只收到一个空公文袋，袋上只注明封套内所装的、可惜实际上并不在内的文件是关于聘任土地测量员一事的。A部呢，一直还在等着我们的回复，这个部虽然已将此事记录在案，但部主管人指望着我们会答复，等着在收到答复之后他们或者聘任土地测量员，或者根据需要继续同我们就此事进行磋商。这种坐等下面回信的事绝非绝无仅有，这是可以理解的，在办什么事都讲求准确、一丝不苟的地方这样做也无可厚非。可是结果呢，因为他指望着我们，所以也就不大重视部里那份有关记录，弄到后来就把这事忘记了。但是在B部就不同，公文袋送到了一位以认真负责闻名的部主管人手里，他叫索尔蒂尼，是个意大利人；连我这个内行人也不明白为什么像他那样能干的人会被长期安置在那个可说很低的职位上。这位索尔蒂尼自然把空文件袋派人送回给我们要求补上内容。但是，从A部发出第一封函件到那时已经过了好几个月，甚至已经好几年了；这也不难理解，因为按常规如果一份文件送对路了，那么最迟在第二天也就送到部里，当天就可以处理；但如果没送对路——在组织非常严密完善的情况下，文件要走这条错误的路简直就必须费很大力气去找，否则是找不到的，那么，那么当然就会拖很长时间了。因此当我们收到索尔蒂尼的便函时，就只能隐隐约约地记起这件事来，当时我们管这事的只有两人，就是米齐和我，那位老师还没有分配给我，我们两个只能对最重要的事保存一份副本，简短地说吧，我们只能给他一个很笼统的回答，说我们完全不知道聘任一事，我们这里也不需要什么土地测量员。”



“但是，”村长顿住了，似乎发觉在乐此不疲的叙述中忘其所以，把话扯远了，或者至少有可能扯远了，“这个故事您听着不觉得腻味吧？”



“不，”K说，“我觉得挺有意思，挺解闷的。”



然后村长说：“我讲这些给您听不是为了给您解闷。”



“我所以觉得它有意思，”K说，“仅仅因为我听了这事后，对那种可笑的混乱算是窥见了一斑，这种混乱有时能决定一个人的命运哟。”



“您还没有真正窥见一斑，”村长正色说，“我可以接着跟您讲。我们的回答当然满足不了索尔蒂尼那样的人。我很佩服这个人，尽管一听到他的名字就头疼。为什么？因为他谁都不信赖，比如，即使他有过无数次机会进行了解，完全知道那是个最可靠的人，可到下一次再碰上时他还是不信任那个人，好像他根本不认识人家似的，或者更正确些说，好像他知道那人是个无赖。我觉得这是对的，一个公职人员就应该这样；可惜我的天性使我不能按这条原则办事，你不是看见了吗，我现在对您这个外乡人把什么底都亮出来了，我这叫做禀性难移，没法子。而索尔蒂尼一接到我们的回信就起了疑。于是便有了一连串的通信往来。索尔蒂尼问我为什么突然想起说不要聘用土地测量员；我仗着米齐的好记性回信说，这事最初是上头提出来的呀（至于事实上是另一个部发来的文件，这一点我们早忘记了）；索尔蒂尼对此的说法是：为什么我到现在才提起上级的这封公函；我又回复他说：因为我现在才想起这封公函来嘛；索尔蒂尼呢：这真是太奇怪了；我：对于拖了那么长时间的一件事，这是一点也不奇怪的；索尔蒂尼：这事确实很奇怪，因为我想起来的那封公函并不存在；我：当然不存在啦，因为关于这事的全部文件都丢失了；索尔蒂尼：如果确有那第一封函件，就必定会有一条有关的记录在，然而这样一条记录并不存在。这时我无话可说了，因为要说是索尔蒂尼那个部出了差错，这话我是既不敢讲也不相信的。您，土地测量员先生，也许会暗暗埋怨索尔蒂尼，心想只要他考虑一下我的话，至少会去向另外的一些部打听打听这事吧。然而真要这样做恰恰是不对的，我不想让您哪怕只是在心里产生一种索尔蒂尼身上还有点毛病的印象。压根不考虑有出错的可能性，正是衙门的一条工作原则。由于组织十分完备严密，这条原则是正确的，而如果想达到极高的办事效率，它也是必要的。所以说索尔蒂尼根本就不能去向其他各部打听，而且就算他去打听，这些部门也决不会答复他，因为它们马上就会感觉出：这是在追究一个出差错的可能性。”



“村长先生，请允许我打断您一下，向您提个问题，”K说，〔8〕“您不是提到过一个检查机关吗？按照您对这里办事情况的描述，如果设想一下这里对这些事居然会没有监督检查，这简直令人不堪忍受。”



“您非常严格，”村长说。“但您就是再严格一千倍，同衙门对待自己的那种严格相比，您的严格也仍然等于零。只有一个地地道道的外乡人才会提出您提的问题。要问有没有检查机关吗？这里有的只是检查机关。当然啰，这些检查机关不是为了查出一般意义上的错误而设立的，因为这里是不会出错的，即便某次出了一个错，如您的情况，但是谁敢下定论，说这肯定是一个错误！”



“这个说法真是太新鲜了！”K叫起来。



“对我来说这是很旧很旧的说法，”村长说。“我同您差不多，也确信是出了一个错，索尔蒂尼因为对这事伤透脑筋毫无进展而得了重病，第一批检查机关官员来了，亏得他们揭出了错误的来源，看出这里有一个错，可是谁敢拍着胸脯说，第二批检查机关官员也会做出同样的判断，第三批和以后各批也都会得出同样的结论？”



“就算是这样吧，”K说，“不过我最好还是不要对这些考虑妄加评论，我也是头一次听说这些检查机关，当然还不可能完全了解它们。但我觉得这里必须区别两点：第一是机关内部发生的事以及仍然从机关的角度出发对这些事可能有这样或那样的看法，第二是我作为一个具体的个人，我这个身在公职机关以外的人，现在正面临着公职机关无视我的存在的危险，这种无视现实是如此荒谬，以致我直到现在总不大相信这危险是实际存在着的。对于第一点，很可能您村长先生刚才讲的那些话都是对的，您对内情了解得如此了如指掌令人叹服，只是我在听了您这番话后也想听您说说我个人的事情了。”



“我也正想来谈这个，”村长说，“但如果我不预先说明几句，那您是听不懂我的话的，我刚才提起检查机关看来是为时过早了。所以现在我要再回过头来谈谈我跟索尔蒂尼的分歧。我已经说过，渐渐的我招架不住了。但索尔蒂尼哪怕只比他的对手略胜一筹，他就等于大获全胜了，因为这样一来他办起事来精力就更集中、精神就更饱满，就更加指挥若定；结果是他的手下败将一见他就胆战心惊，而他手下败将的敌人看着他就会拍案叫绝。只因为我对这后一点在另外一些事情上也有亲身体会，我才能像现在这样讲他这个人。不过我还从未有幸亲眼见过他，他不能到下面来，他真是日理万机，公务繁忙至极，有人对我描述过他的房间，说四壁全被大捆大捆的卷宗挡满，一捆摞一捆，形成了一根根高大的方柱，而且这些还仅仅是索尔蒂尼当时正在处理的文件，由于不断从一捆捆卷宗中抽出、插进文件，并且又都是为赶时间匆匆忙忙地做，这些卷宗堆成的柱子就不断倒塌下来，正是这种接二连三、每隔一会儿就出现一次的轰然巨响，成了索尔蒂尼办公室的突出特征。这也难怪，索尔蒂尼是个一丝不苟的工作人员，不论是最小的事还是最大的事，他都一视同仁，一样认真对待。”



“村长先生，”K说，“您总把我的事叫做小事一桩，但我这件事使许多公职人员忙个不亦乐乎，或许它在开始时只是一桩很小的事吧，可是由于像索尔蒂尼那样的官员们积极参与，它确实已经变成一件大事了。这一点我感到非常遗憾，这是完全违反我的本意的，我不敢有那样的野心：想让关于我的事情的卷宗也变成高大的柱子又随时哗啦啦塌下来，我只希望做个小小的土地测量员，在一张小小的绘图桌旁安安心心地工作就心满意足了。”



“不，”村长说，“您的事不是大事，在这点上您没有理由抱怨，您的事确实是诸多小事中最小的一件。工作量大小并不决定工作的重要性，如果您这样想那您对衙门的了解就太皮毛了。但即使工作量大小能说明问题是否重要，您这件事也仍然是一粒微不足道的芝麻，普遍的事由，即那些没有出所谓差错的事，办理起来工作量就大得多，工作的成果自然也丰富得多，不过您还根本不知道您的事给人造成的实际上的工作量呢，我现在就来说说这个。首先要说的是，索尔蒂尼虽然不再追问我，但他手下的工作人员来了，每天都要在贵宾楼对一些有声望的村民进行几次质询并作翔实记录。大部分人支持我的说法，然而也有少数人起了疑心；土地测量问题触及农民的根本利益，农民们嗅出有人在算计他们，觉得这里头有些不公道的事，他们还找到了一个领头的，于是索尔蒂尼从他们提供的情况中得出结论，确信：如果我当初把问题提到村民大会上讨论，是不会出现全体一致反对聘用土地测量员的情况的。〔9〕这样一来一件原本是明摆着谁都明白的小事——即不需要土地测量员——不管怎么说至少被弄得成为一个问题了。在这件事上，有一个叫布伦斯维克的表现得特别突出——您大概不认识他吧，他也许不是坏人，但又蠢又狂，他是拉塞曼的一个妹夫。”



“是那个鞣皮匠拉塞曼吗？”K问，接着便描述了他在拉塞曼家见到的那个大胡子的长相。



“是的，正是他，”村长说。



“我还认识他的女人，”K有点有一搭没一搭地说。



“这是有可能的，”村长说，然后便沉默了。



“她很美，就是气色不大好，病恹恹的。她大概是城堡里来的人吧？”这后一句话是带着疑问语气说出来的。



村长看了看表，将药水倒进一把汤匙里，急急忙忙灌下肚去。



“您大概只认得城堡里那些办事机构吧？”K颇不客气地问道。



“对，”村长带着自嘲然而却是感激的微笑说道。“机关也是最重要的呵。至于说到布伦斯维克：要是我们能把他从村里驱逐出去，几乎对每个人来说都是件大快人心的事，拉塞曼也不会是最不开心的一个。可是当时布伦斯维克给自己拉了一些支持者：他虽然不是演说家，倒也会大喊大叫，而这样对某些人也能起作用。于是，终归我还是被迫把这件事提交村民大会讨论。不过这也就是布伦斯维克取得的唯一的胜利了，因为到了村民大会上，当然是绝大多数人反对聘用土地测量员的。这次大会也又过去好几年了，但是这些年来这件事一直没有平息，这部分是因为索尔蒂尼太认真，他通过极为周密细致的调查，力图摸清多数派和反对派的动机，再就是因为布伦斯维克的愚蠢和野心，他同官府有各种各样的私人关系，而且有层出不穷的、异想天开的新招数去利用这些关系。当然索尔蒂尼不会上布伦斯维克的当，布伦斯维克怎么瞒得了索尔蒂尼？——可是，恰恰是为了不致上当受骗，需要进行新的调查，新的调查还没有结束，布伦斯维克又想出了新花招，他这人很灵活，这正是他那愚蠢的一种表现。好，现在我来谈谈我们政府机构的一个特殊性质。与它的精细严密相适应，它也极为敏感。如果一件事已经深思熟虑了很久很久，那么即使考虑还没有结束也会出现这种情况：在某个事先无法预见事后也无法找到的地方会闪电般突然冒出一个解决办法，这办法虽说多半是正确地了结了那件事，但无论如何总带有一定的随意性吧。情形就好像是：政府机关再也受不了那绷得太紧的弦、受不了同一件也许本身是微不足道的事情那长年累月的刺激，从而自发地、不等哪个官员来过问就作出决定了。这当然不是说出现了什么奇迹，可以肯定地说是某位官员挥笔写下了这个解决方案，或是作了一个口头的决定，但至少从我们这里，从此地，甚至从公职机关也怎么都不可能弄清楚究竟是哪位官员在这事上作了定夺，是根据什么理由这样决定的了。这些全是检查机关在很久以后才查明的；而我们呢，就永远不得而知了，实际上这事过了那么久恐怕也不会再有谁感兴趣了。我说过，正是这类决定，它们多半是正确的，是很好的，只有一点不足，就是——这类事通常都是如此——一般人知道得太晚，因此在知道前的一段长时间里一直还在对那早已决定了的事进行热烈的讨论。我不知道在您的事情上是否也发出过一个这样的决定。——有迹象说明发出过，也有迹象说明没有发出过；假如已经发出，那么聘任书就可能早已寄给了您，然后您就会长途跋涉到此地来，这些花费了很多时间，而与此同时索尔蒂尼却还会一直在这里为这同一件事操劳着，累得筋疲力尽，布伦斯维克也还会绞尽脑汁搞阴谋诡计，而我就会在他们两人中间受尽了夹板气。这个可能性只是我的揣测，我确切知道的是这一点：一个检查机关查出来好多年前A部曾经就聘任土地测量员一事向村里发出过询问函件，而后直至现在一直没有收到回音。新近又有官员向我查询这事，这样一来全部事实当然也就真相大白了，A部对我提供的不需要土地测量员这一回答表示满意，索尔蒂尼则不得不看到：原来他并不主管这件事情，并且还做了许多无用的、让人伤透脑筋、把人弄得精疲力竭的虚功。——当然，他是无辜的。要不是这段时间也同以往一样新的工作任务不断从四面八方源源而来，要不是您的事情说到底仅仅是桩很小的事——几乎可以说是诸多小事当中最小的一件，那么我们大家一定都会感到如释重负了，我相信索尔蒂尼自己也会这样，只有布伦斯维克有一肚子怨气，但那只会让人觉得可笑罢了。好，现在，恰恰是现在，在整个这件事总算善始善终顺利了结之后——从那时起到今天又已经过去好久，您突然又冒了出来，事情似乎又像要从头开始的样子，请您，土地测量员先生，设想一下我有多么失望呵。现在我下了大决心在我力所能及的范围内决不让这类事再发生，这一点想来您是很清楚也能理解的吧？”
\par\vspace{0.3em}{\footnotesize\color{gray}（第五章）}
\newpage
\section{巴纳巴斯送信}
\vspace{-0.3em}{\color{gray}\small\kaifont \noindent 巴纳巴斯是 K 的信使，也是我们在 Day 2 就认识的人。但直到今天，他的姐姐奥尔嘉才把他每天在城堡办公厅里经历的事情讲清楚。他没有工作服，所以他连自己算不算城堡的人都不知道。他穿梭在栅栏之间，那些允许他穿过的，和不允许他穿过的，外表一模一样。克拉姆给他派任务——听起来很荣耀——但信是秘书写的，内容是过时的。办公厅里的设置不是为了让人找到正确的路，是为了让外面的人找不到任何路。}\vspace{0.5em}
\vspace{0.3em}

\noindent\rule{\textwidth}{0.4pt}

\vspace{0.3em}



总而言之，巴纳巴斯绝对谈不上是职位较高的办事人员，那么照理说他就应该是个低级的办事人员了吧？可是那些低级办事人员全都有工作服，至少，在他们下到村里来时是穿着工作服的，这不是真正的制服，各人穿的也有很多小地方不一样，可是不管怎么说看到这种服装一眼便能认出是城堡来的人员，在贵宾楼你不是也见到过这样一些人嘛。这种服装最引人注目的一点就是多半都非常紧身，干农活的或者做手工的人恐怕是没法穿的。好了，巴纳巴斯没有这种衣服；这不光是让人感觉寒碜，觉得脸上无光，这倒还能忍受，严重的是让你对什么事都怀疑起来，特别是在情绪低落消沉时更是这样——我们，就是巴纳巴斯和我，我们有时会情绪很低落的，我说有时，并不是说这样的时候就很少很少。遇上这种时候我们就对什么都怀疑起来，我们琢磨：巴纳巴斯做的工作到底算不算是城堡的工作；不错，他是在那些办公厅里进进出出，但那些办公厅就是真正的城堡吗？就说城堡肯定有不少办公厅吧，但是，允许巴纳巴斯进出的那些办公厅是城堡的办公厅吗？他是走进了一些办公厅；可那毕竟只是所有办公厅中的一部分呀，那儿另有些栅栏门，栅栏门后面还有另外一些办公厅。人家倒也没有禁止他穿过这些栅栏，可他不可能再往前走了，因为在栅栏这边他就找到了他的那几个上司，他们听完他的报告就让他走人了。另外，在那里行动一直是受到监视的，至少到过那里的人有这个感觉。再说即使他穿过了栅栏又有什么用处呢，他在那边又没有要办的事，那样做岂不是成了一个擅自闯入公务重地的异己分子？你也不能把这些栅栏设想成某种不许越过的界线，巴纳巴斯也一再让我注意这一点。原来，在他进出的那些办公厅里也是有栅栏的；这就是说也有那么一些栅栏是他经常穿行的，而且这些栅栏跟他还没有穿越的那些看上去并没有两样，所以说也不能一开始就猜想在他没有穿越的栅栏后面会有跟他去过的那些办公厅大不相同的什么别的办公厅。不就是我们只在情绪低落消沉时才有这种想法嘛！就这样，我们的疑心越来越大，根本无法抗拒。再举个例子说吧，巴纳巴斯是跟一些官员说话，巴纳巴斯是得到人家让他传送的信，可那究竟是些什么官员，是些什么信呵！他自己说他现在是分配到克拉姆手下了，克拉姆本人亲自给他派任务。好，这应该说是很抬举他了吧，连一些高级管事也没有争取到这份差事，唔，恐怕这是把他抬得太高了，真让人担心会摔个粉身碎骨的。你想想看，直接分给克拉姆调遣，同他嘴对嘴地说话！真是这么回事吗？



...



你知道，弗丽达不怎么喜欢我，她是怎么也不乐意让我见到克拉姆的，但他的相貌在村里当然是人人都知道的，有少数人见过他，所有的人都听说过他，于是从见过的人得到的印象、道听途说听来的材料、再加上某些别有用心的人的添油加醋的描述，就勾画出了一幅克拉姆画像，这画像也许大体上符合实际，但只能说是大概像。除基本轮廓外，别的方面就有比较大的出入了，而这些描述的出入之处也许还没有克拉姆的真正外貌的变化来得大呢。据说他到村子里来时完全变了模样，而到离开村子时又变一个样，喝啤酒前是一个样，喝完后又是另一个样，醒时一个样，睡着了又一个样，独自一人时一个样，跟人谈话时又一个样，根据这些，说他在上头城堡里时完完全全是另一个人就由不得你不信了。再就是，即使那些有关他在村里时的外貌的说法，也有相当大的出入，关于他的个子、姿势、胖瘦、胡子等方面的说法都很不一致，只有在穿着一点上，幸好大家的说法总算一致了：都说他总是穿同一件衣服，一件有长摆的黑色上装。为什么会有这么些不一样的说法？当然不是克拉姆会变魔术，而是因为见到他的人当时的心情、激动的程度、各自抱的希望或绝望的复杂心境千差万别，数也数不清，加上见到他的人多半都只被允许匆匆一面，考虑到这些因素，有那么多不同说法就非常容易理解了。
\par\vspace{0.3em}{\footnotesize\color{gray}（第十五章）}
\newpage
\begin{center}
{\LARGE \bfseries Day 7\\[1em]收束与回望}
\end{center}
\vspace{3em}
\section{聘用始末}
\vspace{-0.3em}{\color{gray}\small\kaifont \noindent K 不只要一个简单的说法。他想搞清楚，聘他来做土地测量员的这个决定是如何做出的。村长讲了很久。故事的核心不是答案。}\vspace{0.5em}
\vspace{0.3em}

\noindent\rule{\textwidth}{0.4pt}

\vspace{0.3em}



“这么说这两个助手，”村长又开口了，洋洋得意地微笑着，似乎在说：所有这些全是他一手安排的，而谁都无法预想到这一点，“这么说，他们让您感到讨厌，可他们是您自己的助手呵。”——“不对，”K冷冷地说，“他们是我到此地后才跑到我身边来的。”——“怎么，跑来的，”村长说，“您的意思大概是说派来的吧。”——“好吧，就算是派来的吧，”K说。“但说他们是从天外飞来的也未尝不可，这种派法也够欠考虑的了。”——“这里办什么事都不会欠考虑，”村长说，急得连腿疼也忘了，一骨碌坐起来。“办什么事都不欠考虑，”K说，“那么招聘我这事又怎么说？”——“聘用您也是经过深思熟虑的，”村长说，“只是有一些次要因素起了干扰作用罢了，我可以用有关的文件向您证明这一点。”——“这些文件恐怕找不到了，”K说。“找不到？”村长叫起来。“米齐，请快点找！不过，没有文件我也可以先把事情的经过讲给您听听。刚才我提到的那份指令，我们接到后当即回复上头表示感谢，说我们不需要土地测量员。但是这封回信看来是没有送回发出指令的那个部去，我管它叫A部吧，而是误送到另一个部，送到B部去了。这样A部就一直没有得到回音，然而可惜连B部也没有收到我们复信的全部；一种可能是公文袋里装的信落在村公所，再一种可能是在半路上遗失了——肯定不会在部里遗失，这我敢担保，总之，B部也只收到一个空公文袋，袋上只注明封套内所装的、可惜实际上并不在内的文件是关于聘任土地测量员一事的。A部呢，一直还在等着我们的回复，这个部虽然已将此事记录在案，但部主管人指望着我们会答复，等着在收到答复之后他们或者聘任土地测量员，或者根据需要继续同我们就此事进行磋商。这种坐等下面回信的事绝非绝无仅有，这是可以理解的，在办什么事都讲求准确、一丝不苟的地方这样做也无可厚非。可是结果呢，因为他指望着我们，所以也就不大重视部里那份有关记录，弄到后来就把这事忘记了。但是在B部就不同，公文袋送到了一位以认真负责闻名的部主管人手里，他叫索尔蒂尼，是个意大利人；连我这个内行人也不明白为什么像他那样能干的人会被长期安置在那个可说很低的职位上。这位索尔蒂尼自然把空文件袋派人送回给我们要求补上内容。但是，从A部发出第一封函件到那时已经过了好几个月，甚至已经好几年了；这也不难理解，因为按常规如果一份文件送对路了，那么最迟在第二天也就送到部里，当天就可以处理；但如果没送对路——在组织非常严密完善的情况下，文件要走这条错误的路简直就必须费很大力气去找，否则是找不到的，那么，那么当然就会拖很长时间了。因此当我们收到索尔蒂尼的便函时，就只能隐隐约约地记起这件事来，当时我们管这事的只有两人，就是米齐和我，那位老师还没有分配给我，我们两个只能对最重要的事保存一份副本，简短地说吧，我们只能给他一个很笼统的回答，说我们完全不知道聘任一事，我们这里也不需要什么土地测量员。”



“但是，”村长顿住了，似乎发觉在乐此不疲的叙述中忘其所以，把话扯远了，或者至少有可能扯远了，“这个故事您听着不觉得腻味吧？”



“不，”K说，“我觉得挺有意思，挺解闷的。”



然后村长说：“我讲这些给您听不是为了给您解闷。”



“我所以觉得它有意思，”K说，“仅仅因为我听了这事后，对那种可笑的混乱算是窥见了一斑，这种混乱有时能决定一个人的命运哟。”



“您还没有真正窥见一斑，”村长正色说，“我可以接着跟您讲。我们的回答当然满足不了索尔蒂尼那样的人。我很佩服这个人，尽管一听到他的名字就头疼。为什么？因为他谁都不信赖，比如，即使他有过无数次机会进行了解，完全知道那是个最可靠的人，可到下一次再碰上时他还是不信任那个人，好像他根本不认识人家似的，或者更正确些说，好像他知道那人是个无赖。我觉得这是对的，一个公职人员就应该这样；可惜我的天性使我不能按这条原则办事，你不是看见了吗，我现在对您这个外乡人把什么底都亮出来了，我这叫做禀性难移，没法子。而索尔蒂尼一接到我们的回信就起了疑。于是便有了一连串的通信往来。索尔蒂尼问我为什么突然想起说不要聘用土地测量员；我仗着米齐的好记性回信说，这事最初是上头提出来的呀（至于事实上是另一个部发来的文件，这一点我们早忘记了）；索尔蒂尼对此的说法是：为什么我到现在才提起上级的这封公函；我又回复他说：因为我现在才想起这封公函来嘛；索尔蒂尼呢：这真是太奇怪了；我：对于拖了那么长时间的一件事，这是一点也不奇怪的；索尔蒂尼：这事确实很奇怪，因为我想起来的那封公函并不存在；我：当然不存在啦，因为关于这事的全部文件都丢失了；索尔蒂尼：如果确有那第一封函件，就必定会有一条有关的记录在，然而这样一条记录并不存在。这时我无话可说了，因为要说是索尔蒂尼那个部出了差错，这话我是既不敢讲也不相信的。您，土地测量员先生，也许会暗暗埋怨索尔蒂尼，心想只要他考虑一下我的话，至少会去向另外的一些部打听打听这事吧。然而真要这样做恰恰是不对的，我不想让您哪怕只是在心里产生一种索尔蒂尼身上还有点毛病的印象。压根不考虑有出错的可能性，正是衙门的一条工作原则。由于组织十分完备严密，这条原则是正确的，而如果想达到极高的办事效率，它也是必要的。所以说索尔蒂尼根本就不能去向其他各部打听，而且就算他去打听，这些部门也决不会答复他，因为它们马上就会感觉出：这是在追究一个出差错的可能性。”



“村长先生，请允许我打断您一下，向您提个问题，”K说，〔8〕“您不是提到过一个检查机关吗？按照您对这里办事情况的描述，如果设想一下这里对这些事居然会没有监督检查，这简直令人不堪忍受。”



“您非常严格，”村长说。“但您就是再严格一千倍，同衙门对待自己的那种严格相比，您的严格也仍然等于零。只有一个地地道道的外乡人才会提出您提的问题。要问有没有检查机关吗？这里有的只是检查机关。当然啰，这些检查机关不是为了查出一般意义上的错误而设立的，因为这里是不会出错的，即便某次出了一个错，如您的情况，但是谁敢下定论，说这肯定是一个错误！”



“这个说法真是太新鲜了！”K叫起来。



“对我来说这是很旧很旧的说法，”村长说。“我同您差不多，也确信是出了一个错，索尔蒂尼因为对这事伤透脑筋毫无进展而得了重病，第一批检查机关官员来了，亏得他们揭出了错误的来源，看出这里有一个错，可是谁敢拍着胸脯说，第二批检查机关官员也会做出同样的判断，第三批和以后各批也都会得出同样的结论？”



“就算是这样吧，”K说，“不过我最好还是不要对这些考虑妄加评论，我也是头一次听说这些检查机关，当然还不可能完全了解它们。但我觉得这里必须区别两点：第一是机关内部发生的事以及仍然从机关的角度出发对这些事可能有这样或那样的看法，第二是我作为一个具体的个人，我这个身在公职机关以外的人，现在正面临着公职机关无视我的存在的危险，这种无视现实是如此荒谬，以致我直到现在总不大相信这危险是实际存在着的。对于第一点，很可能您村长先生刚才讲的那些话都是对的，您对内情了解得如此了如指掌令人叹服，只是我在听了您这番话后也想听您说说我个人的事情了。”



“我也正想来谈这个，”村长说，“但如果我不预先说明几句，那您是听不懂我的话的，我刚才提起检查机关看来是为时过早了。所以现在我要再回过头来谈谈我跟索尔蒂尼的分歧。我已经说过，渐渐的我招架不住了。但索尔蒂尼哪怕只比他的对手略胜一筹，他就等于大获全胜了，因为这样一来他办起事来精力就更集中、精神就更饱满，就更加指挥若定；结果是他的手下败将一见他就胆战心惊，而他手下败将的敌人看着他就会拍案叫绝。只因为我对这后一点在另外一些事情上也有亲身体会，我才能像现在这样讲他这个人。不过我还从未有幸亲眼见过他，他不能到下面来，他真是日理万机，公务繁忙至极，有人对我描述过他的房间，说四壁全被大捆大捆的卷宗挡满，一捆摞一捆，形成了一根根高大的方柱，而且这些还仅仅是索尔蒂尼当时正在处理的文件，由于不断从一捆捆卷宗中抽出、插进文件，并且又都是为赶时间匆匆忙忙地做，这些卷宗堆成的柱子就不断倒塌下来，正是这种接二连三、每隔一会儿就出现一次的轰然巨响，成了索尔蒂尼办公室的突出特征。这也难怪，索尔蒂尼是个一丝不苟的工作人员，不论是最小的事还是最大的事，他都一视同仁，一样认真对待。”



“村长先生，”K说，“您总把我的事叫做小事一桩，但我这件事使许多公职人员忙个不亦乐乎，或许它在开始时只是一桩很小的事吧，可是由于像索尔蒂尼那样的官员们积极参与，它确实已经变成一件大事了。这一点我感到非常遗憾，这是完全违反我的本意的，我不敢有那样的野心：想让关于我的事情的卷宗也变成高大的柱子又随时哗啦啦塌下来，我只希望做个小小的土地测量员，在一张小小的绘图桌旁安安心心地工作就心满意足了。”



“不，”村长说，“您的事不是大事，在这点上您没有理由抱怨，您的事确实是诸多小事中最小的一件。工作量大小并不决定工作的重要性，如果您这样想那您对衙门的了解就太皮毛了。但即使工作量大小能说明问题是否重要，您这件事也仍然是一粒微不足道的芝麻，普遍的事由，即那些没有出所谓差错的事，办理起来工作量就大得多，工作的成果自然也丰富得多，不过您还根本不知道您的事给人造成的实际上的工作量呢，我现在就来说说这个。首先要说的是，索尔蒂尼虽然不再追问我，但他手下的工作人员来了，每天都要在贵宾楼对一些有声望的村民进行几次质询并作翔实记录。大部分人支持我的说法，然而也有少数人起了疑心；土地测量问题触及农民的根本利益，农民们嗅出有人在算计他们，觉得这里头有些不公道的事，他们还找到了一个领头的，于是索尔蒂尼从他们提供的情况中得出结论，确信：如果我当初把问题提到村民大会上讨论，是不会出现全体一致反对聘用土地测量员的情况的。〔9〕这样一来一件原本是明摆着谁都明白的小事——即不需要土地测量员——不管怎么说至少被弄得成为一个问题了。在这件事上，有一个叫布伦斯维克的表现得特别突出——您大概不认识他吧，他也许不是坏人，但又蠢又狂，他是拉塞曼的一个妹夫。”



“是那个鞣皮匠拉塞曼吗？”K问，接着便描述了他在拉塞曼家见到的那个大胡子的长相。



“是的，正是他，”村长说。



“我还认识他的女人，”K有点有一搭没一搭地说。



“这是有可能的，”村长说，然后便沉默了。



“她很美，就是气色不大好，病恹恹的。她大概是城堡里来的人吧？”这后一句话是带着疑问语气说出来的。



村长看了看表，将药水倒进一把汤匙里，急急忙忙灌下肚去。



“您大概只认得城堡里那些办事机构吧？”K颇不客气地问道。



“对，”村长带着自嘲然而却是感激的微笑说道。“机关也是最重要的呵。至于说到布伦斯维克：要是我们能把他从村里驱逐出去，几乎对每个人来说都是件大快人心的事，拉塞曼也不会是最不开心的一个。可是当时布伦斯维克给自己拉了一些支持者：他虽然不是演说家，倒也会大喊大叫，而这样对某些人也能起作用。于是，终归我还是被迫把这件事提交村民大会讨论。不过这也就是布伦斯维克取得的唯一的胜利了，因为到了村民大会上，当然是绝大多数人反对聘用土地测量员的。这次大会也又过去好几年了，但是这些年来这件事一直没有平息，这部分是因为索尔蒂尼太认真，他通过极为周密细致的调查，力图摸清多数派和反对派的动机，再就是因为布伦斯维克的愚蠢和野心，他同官府有各种各样的私人关系，而且有层出不穷的、异想天开的新招数去利用这些关系。当然索尔蒂尼不会上布伦斯维克的当，布伦斯维克怎么瞒得了索尔蒂尼？——可是，恰恰是为了不致上当受骗，需要进行新的调查，新的调查还没有结束，布伦斯维克又想出了新花招，他这人很灵活，这正是他那愚蠢的一种表现。好，现在我来谈谈我们政府机构的一个特殊性质。与它的精细严密相适应，它也极为敏感。如果一件事已经深思熟虑了很久很久，那么即使考虑还没有结束也会出现这种情况：在某个事先无法预见事后也无法找到的地方会闪电般突然冒出一个解决办法，这办法虽说多半是正确地了结了那件事，但无论如何总带有一定的随意性吧。情形就好像是：政府机关再也受不了那绷得太紧的弦、受不了同一件也许本身是微不足道的事情那长年累月的刺激，从而自发地、不等哪个官员来过问就作出决定了。这当然不是说出现了什么奇迹，可以肯定地说是某位官员挥笔写下了这个解决方案，或是作了一个口头的决定，但至少从我们这里，从此地，甚至从公职机关也怎么都不可能弄清楚究竟是哪位官员在这事上作了定夺，是根据什么理由这样决定的了。这些全是检查机关在很久以后才查明的；而我们呢，就永远不得而知了，实际上这事过了那么久恐怕也不会再有谁感兴趣了。我说过，正是这类决定，它们多半是正确的，是很好的，只有一点不足，就是——这类事通常都是如此——一般人知道得太晚，因此在知道前的一段长时间里一直还在对那早已决定了的事进行热烈的讨论。我不知道在您的事情上是否也发出过一个这样的决定。——有迹象说明发出过，也有迹象说明没有发出过；假如已经发出，那么聘任书就可能早已寄给了您，然后您就会长途跋涉到此地来，这些花费了很多时间，而与此同时索尔蒂尼却还会一直在这里为这同一件事操劳着，累得筋疲力尽，布伦斯维克也还会绞尽脑汁搞阴谋诡计，而我就会在他们两人中间受尽了夹板气。这个可能性只是我的揣测，我确切知道的是这一点：一个检查机关查出来好多年前A部曾经就聘任土地测量员一事向村里发出过询问函件，而后直至现在一直没有收到回音。新近又有官员向我查询这事，这样一来全部事实当然也就真相大白了，A部对我提供的不需要土地测量员这一回答表示满意，索尔蒂尼则不得不看到：原来他并不主管这件事情，并且还做了许多无用的、让人伤透脑筋、把人弄得精疲力竭的虚功。——当然，他是无辜的。要不是这段时间也同以往一样新的工作任务不断从四面八方源源而来，要不是您的事情说到底仅仅是桩很小的事——几乎可以说是诸多小事当中最小的一件，那么我们大家一定都会感到如释重负了，我相信索尔蒂尼自己也会这样，只有布伦斯维克有一肚子怨气，但那只会让人觉得可笑罢了。好，现在，恰恰是现在，在整个这件事总算善始善终顺利了结之后——从那时起到今天又已经过去好久，您突然又冒了出来，事情似乎又像要从头开始的样子，请您，土地测量员先生，设想一下我有多么失望呵。现在我下了大决心在我力所能及的范围内决不让这类事再发生，这一点想来您是很清楚也能理解的吧？”
\par\vspace{0.3em}{\footnotesize\color{gray}（第五章）}
\newpage
\section{根茎}
\vspace{-0.3em}{\color{gray}\small\kaifont \noindent 德勒兹与加塔利说，世界上有三种生命空间形态。树根，统一的生命体，从唯一起点出发，通过二元分化不断生长。胚根，矛盾的生命体，多个起点相互嫁接，任何路径都可以被替代，整体在崩溃和重建之间轮回。根茎，多元的生命体，脱离起源，不断迁移、断裂，在任何一个节点上都蕴含着重新开始的潜力。}\vspace{0.5em}
\vspace{0.3em}

\noindent\rule{\textwidth}{0.4pt}

\vspace{0.3em}



书的第一种类型，是“书—根”（livre-racine）。树已经是世界的形象，或者说，根是“世界—树”的形象。这是经典之书，作为有机的、意谓的、主观的崇高内在性（书之尽）。书模仿世界，正如艺术模仿自然；运用其自身所特有的手段，产生出那些自然不能或不再能创造出来的东西。书的法则，则是反思／映像（réflexion）的法则，即“一”生成为“二”。那书的法则又怎能存在于自然之中，既然它统辖着世界与书、自然与艺术之间的区分？一生成二；每当我们遇到这个法则——即便被毛泽东策略性地提出，或被人们最为“辩证地”理解，我们所面临着的正是最为经典、最具反思性、最古老、最令人厌倦的思想。自然并不是这样运作的；在自然之中，根是直根，具有更为大量的侧面的或环状的（而非二元分化的）分支。精神落后于自然。甚至作为自然实在的书也是一种直根，有着垂直的轴和环绕的叶。然而书作为精神实在——树或根的形象——却不断展现着一生二、二生四的法则……二元逻辑是“树—根”的精神实在。即使像语言学这样“领先的”学科也仍然恪守着“树—根”这个基本形象，由此与反思的古典传统联结在一起（比如乔姆斯基和他的意群树：从一个点S开始，以二元分化的方式衍生）。可以说，此种思想从未理解多元性：它必需预设一种根本的、强有力的统一，以便遵循一种精神的方法来达到“二”。从客体的角度来说，遵循自然的方法，人们无疑可以直接从“一”达到三、四或五，但却始终必须预设一种根本的、强有力的统一性作为前提，即支撑着次级根的直根的统一性。这并不能使我们走得更远。取代二元分化之二元逻辑的，只不过是连续的循环之间的——对应的关系。无论是直根还是二歧根，都不能更好地解释多元性。前者运作于客体之中，后者运作于主体之中。二元逻辑和一一对应关系仍然统治着精神分析［弗洛伊德在解释施列伯（Schreber）病例时所运用的谵妄之树］、语言学、结构主义，乃至信息论。



胚根系统，或束根，是书的第二种形象，我们的现代性显然仰赖于它。这回，主根已然夭折，或者，它的末端已然被摧毁，一个任意的、直接的次级根的多元体被嫁接于它之上，并呈现出一种蓬勃的生长态势。这回，自然实在显现于主根夭折的过程中，但根的统一性仍然持有，作为过去或将来，作为可能性。我们须追问，是否反思性的精神实在没有以这样的方式对此种事态进行弥补：即转而诉求一种更为全面的隐秘统一性或更为广泛的总体性。不妨采用威廉·巴勒斯的剪裁法（cut-up）：将一篇文本叠合（pliage）进另一篇之中，由此构成了多元的，甚至是偶生的根（比如一根插穗），这对于所考察的文本来说就意味着一种替补的维度（dimension supplémentaire）。正是在这个叠合的替补维度之中，统一性继续着其精神的劳作。正是在这个意义上，最为碎片化的著作也同样可以被视为“全集”或“巨著”。绝大多数用来令系列得以增殖或令多元性得以拓展的现代方法在某个方向上（比如线性的方向）是极为有效的，而一种总体化的统一性却在另一个维度（循环或周期）之中获得了更为有力的肯定。每当一个多元体被掌控于一种结构之中时，其增长就被后者的组合法则所缩减甚至抵消。主张摧毁统一性的人也正是天使的缔造者，**天使博士（*doctores angelici*）\textbf{，因为他们肯定了一种真正的、至上的、天使般的统一性。乔伊斯的词语可谓名副其实地“具有多重根”，它们有效地粉碎了词语乃至语言的线性统一，但由此却得出了一种句子、文本或知识的循环统一。尼采的格言粉碎了知识的线性统一，但却由此回归于永恒轮回的循环统一［作为思想之中的非—知（non-su）］。如此说来，束根系统尚未真正摆脱二元论，它展现出了主客之间、自然实在和精神实在之间的互补性：在客体之中，统一性不断遭到阻碍甚至阻止，但在主体之中，一种新的统一性却获得胜利。世界失去了其主根，而主体也不再能够进行二元分化，但却在于其客体的维度始终构成替补的维度中，达到了一种更高的统一性，一种矛盾的或超定（surdétermination）的统一性。世界生成为混沌，但书仍然是世界的形象；胚根—混沌体（chaosmos），而不再是根—宇宙（cosmos）。说异的迷思：因其碎片化而更具有总体性的书。书作为世界的形象，无论如何都是乏味的观念。确实，仅仅说“多元性万岁”，这还是不够的，尽管要想发出这样的呼声也绝非易事。没有哪种排版、词语、甚或句法的技巧足以使其被听见。}必须形成“多”，但不是通过始终增加一个更高的维度，而是相反，以最为简单而节制的方式，始终对已有的维度进行 n-1 的操作（只有如此，一才成为多的构成部分，即始终是被减去）。从有待构成的多元体中减去独一无二者，在 n-1 的维度上写作。这样的体系可以被称为根茎（rhizome）。作为地下的茎，根茎截然有别于根和胚根。球茎和块茎都属于根茎。具有根或胚根的植物从所有其他的方面来看也可以是根茎式的；问题就在于，是否植物学从其特性上来说是完全是根茎式的。甚至某些动物也是根茎式的，在其集群（meute）的形态之中。鼠群就是根茎。兽穴也是根茎，它们有着各种栖居、储藏、移动、躲避、断裂的功能。根茎自身具有异常多样的形态，从表面向四面八方延展，直到凝聚成球茎和块茎的形态。当鼠群之间彼此窜动之时。存在着最好的和最糟的根茎：土豆，茅草和莠草。动物和植物，茅草就是螃蟹草（crab-grass）。我们觉得，如果不列举一些根茎的大致特征的话，是无法令人信服的。

\newpage
\vspace*{1em}
\begin{center}
{\Large\bfseries 课程综述}
\end{center}
\vspace{1em}
\addcontentsline{toc}{section}{课程综述}

从哪里开始，又走到了哪里。



\textbf{第一天。} 你们变成了一支田野调查队的成员。不是来评论一本小说，而是潜入陌生的权力场域。K 已经在村子里了，给你们写了第一封情报。你们接触了田野伦理，读了人类学调查方法的清单。你们意识到，田野调查是一种既疏离又热烈的思想姿态。观察与干预相伴而行，任何一次策略性抉择，都意味着真实之树的一次分叉。你们向城堡投去了最初的一瞥。



\textbf{第二天。} K 泊入田野的行动相当激进，旋即与田野中人建立亲密关系。这引发了你们的担忧和争论，是深入田野的重大进展，还是一步险棋。你们注意到 K 把同一个名字冠在两个老助手头上，一个叫阿图尔一个叫耶利米，但他们都叫“助手”。这是被动之举还是防范措施？你们感到不适，心怀求知的渴望却迷失于细节和疑团里，是煎熬的。但你们意识到，这正是调查的常态，关键在于自觉地意识到它。想要解释职能与名称的分离，就要理解政治分工有别于经济分工的奥秘。



\textbf{第三天。} K 完成了田野角色的确立。你们追问，远离权力中心，为什么反而是介入组织的更好途径。成为组织的一员，将对个体产生哪些心理后果。研究城堡、组织、权力，将如何有助于改善我们的社会生活。你们积累了很多问题，准备在后面的时间里去深化和解决它们。



\textbf{第四天。} 你们读了周雪光，了解了制度化了的政府间共谋现象。为什么组织权威默许共谋？讨论里出现了不同的声音，有人认为是有心无力鞭长莫及，有人认为相安无事就是更好的选择。你们理解了决策-执行的二元关系，执行对决策的扭曲不可避免，决策者必须在决策环节预判和吸纳这种扭曲。这才是非正式制度的真正来源。但你们没有停在这里，如果非正式的执行行为已经制度化，为什么不直接正式化？问题背后是体制之辩，是治理空间结构的潜在多元性。



\textbf{第五天。} 你们和纪君班合班上课，经历了一场从现象到理论的思想历程。你们分享了对教师角色的理解，意识到教师身处复杂的关系网络中，在制度之下努力维系着诸多关乎个人的良好意愿。但个体的良好意愿不易上升为组织行为。组织对确定性的偏爱促使教师倡导标准和可预测，哪怕他们把学生的差异和潜力看在心里。老师是在学生生活中栖居的街头官僚，正式职责之外，固有着一系列由栖居于他人生活而带来的非正式职责。两种职责并不平衡，这一定程度上源于家长特别的中介地位。你们也读了 1985 年的教育体制改革决定，教育被设定为服务于经济和行政建制的人才供应线。你们有理由相信，这样富有时代色彩的选择将随时间而改变。



\textbf{第六天。} 你们回到了城堡本身。不是在外部看，而是进入它的内部运作。电话是八音盒，打到城堡的电话让最下一级部门全部响起来，但接起来的人不一定是你找的人。公文在每个环节都有回应，但从未有裁断。下级负责遗忘，把决定权推出去，上级负责铭记，不断追问、催促。这组织的动态平衡。



\textbf{第七天。} 德勒兹与加塔利说，世界上有三种生命空间形态。树根，统一的生命体，从唯一起点出发，通过二元分化不断生长，对应的是韦伯式的正式组织。胚根，矛盾的生命体，多个起点相互嫁接，任何路径都可以被替代，整体在崩溃和重建之间轮回，对应的是列宁式的革命组织。根茎，多元的生命体，脱离起源，不断迁移、断裂，在任何一个节点上都蕴含着重新开始的潜力，对应的是卡夫卡式的非正式组织。



三种形态是理想型。真实世界中的组织从未只是其中一个。它们声称为韦伯式的正式组织，渴望成为列宁式的革命组织，而实际上是卡夫卡式的非正式组织。后者如幽灵般弥漫，难以捉摸，又如真相般难以抗拒。



这就是你们走过的七天。从田野调查者到组织观察者，从读一本书到理解一种存在方式。七天前你们打开《城堡》，可能觉得它是一个关于官僚制的寓言。现在你们知道，它不是寓言。它是田野笔记。



每个人看到的城堡都不一样。不是因为有人看对了有人看错了，而是因为理解组织就是理解自己理解它的方式。你们意识到这一点，这七天就值得了。



谢谢你们的思想劳作，感谢这一程的陪伴！
\end{document}
